\reviewersec{Part 2: Response to Reviewer 2}






\subsection*{R2-1: Scientific Innovation and Novelty}
\revcomment{R2-1}{> High scientific novelty: The integration of 3D rectified flow, a domain-adapted BiomedCLIP/VLM, and multi-tracer PET synthesis is a clear methodological contribution beyond conventional MRI-to-PET translation.}

\response{
We sincerely thank the reviewer for the positive assessment of the methodological contribution of our work. We appreciate the recognition of the value of integrating 3D rectified flow, domain-adapted BiomedCLIP/VLM, and multi-tracer PET synthesis.
}


\actiontaken{
We have carefully verified this point. No further manuscript text modifications were required beyond clarifying the discussion.
}


\subsection*{R2-2: Clinical Motivation and Radiation Reduction}
\revcomment{R2-2}{> Strong clinical motivation: Synthesizing AV45-PET and FDG-PET from routine MRI addresses important limitations related to PET cost, radiation exposure, and accessibility.}

\response{
We thank the reviewer for highlighting the clinical relevance of our study. Addressing the challenges of PET acquisition costs, radiation exposure, and limited accessibility through MRI-based synthesis of AV45-PET and FDG-PET is a central motivation for our work.
}


\actiontaken{
We have carefully verified this point. No further manuscript text modifications were required beyond clarifying the discussion.
}


\subsection*{R2-3: Cross-Validation Protocols for Synthesis vs. Classification}
\revcomment{R2-3}{> Cross-validation inconsistency: The Results section reports 3-fold cross-validation, whereas the preprocessing/methodology section reports a two-fold scheme. The exact protocol must be reconciled and stated consistently for each experiment.}

\response{
We thank the reviewer for pointing out the insufficient clarification of the cross-validation protocols in the original manuscript.

We confirm that the two-fold and three-fold cross-validation schemes were employed for different experimental stages. Specifically, two-fold cross-validation was used for PET image synthesis, where the training and held-out subsets were exchanged between folds to generate out-of-fold (OOF) synthetic PET images for the study subjects. Subsequently, a separate three-fold cross-validation scheme was applied to the downstream diagnostic classification tasks using the generated synthetic PET images to evaluate the contribution of different imaging inputs to classification performance.

Therefore, the two cross-validation schemes are not contradictory but correspond to distinct evaluation stages: PET image synthesis and downstream diagnostic classification. We will clarify their respective purposes and consistently distinguish these protocols in the Results and Methods sections of the revised manuscript to avoid ambiguity.
}

% [Audit Note: 上述历史三折描述与主页面“原始分类三折重复使用同一折”的已确认事实存在冲突；新实验完成前不得将其视为历史有效三折结果。]

\actiontaken{
We have revised the manuscript accordingly across the following sections:
\begin{itemize}[leftmargin=1.5em, itemsep=2pt]
  \item \textbf{Results: Superior Fidelity and Generalizability of Synthesized PET, Line 102}
  \item \textbf{Methods: Materials and Data Preprocessing, Line 303}
\end{itemize}
}

\changedtext{Results: Superior Fidelity and Generalizability of Synthesized PET, Line 102}{%
Models were trained on the ADNI dataset and evaluated under two settings: internally via two-fold cross-validation for PET synthesis on ADNI, and externally on the independent OASIS cohort (for which only $^{18}$F-AV-45 was available).%
}

% [Revision Note: 将错误的 `3-fold` 修正为 `two-fold`，同时增加 `for PET synthesis` 明确适用范围。其余内容保持不变。]
% [Revision Note: 修改二：Methods → Materials and Data Preprocessing]
% [Revision Note: 位置：`BMEF_main.tex` 第 303 行。]
\changedtext{Methods: Materials and Data Preprocessing, Line 303}{%
PET synthesis was evaluated under a two-fold cross-validation scheme: the cohort was randomly partitioned into two equal-sized subsets, with alternating roles as training/validation and test sets to ensure unbiased performance estimation. The held-out synthetic PET images generated from both folds were combined to form an out-of-fold (OOF) dataset for subsequent analyses.%
}

% [Revision Note: 将 `Model evaluation` 明确为 `PET synthesis`，并补充一句 OOF 合并与下游分类之间的关系。]
% [Revision Note: 这与原作者确认的历史实验流程一致，也解释了为什么生成阶段采用两折，而分类阶段可以独立采用三折。]

\subsection*{R2-4: Cross-Sectional Staging vs. Longitudinal Prognostic Claims}
\revcomment{R2-4}{> Overstatement of prognostic stratification: EMCI-versus-LMCI classification is a stage-discrimination task and does not by itself establish longitudinal prognosis or future conversion to AD. The terminology should be moderated unless longitudinal conversion outcomes are analyzed.}

\response{
We thank the reviewer for highlighting the important distinction between MCI stage discrimination and longitudinal prognostic prediction.

Our original motivation for relating EMCI/LMCI classification to disease progression was based on previous longitudinal studies showing that patients with LMCI generally exhibit more severe cognitive impairment and have a higher subsequent risk of conversion to dementia than those with EMCI (Aisen et al., 2015; Lin et al., 2022; Petersen, 2021). This clinical evidence motivated us to include EMCI-versus-LMCI classification as a downstream task to evaluate the performance of synthetic PET in distinguishing clinically relevant MCI stages.

Nevertheless, we agree that such population-level associations do not establish individual-level prognostic prediction. Our study evaluated cross-sectional classification performance without analyzing longitudinal conversion outcomes. Therefore, the EMCI-versus-LMCI classification results cannot, by themselves, demonstrate the ability to predict future conversion to AD or disease progression.

To address this concern, we will explicitly describe this task as MCI stage discrimination and revise the relevant terminology and conclusions throughout the manuscript. In particular, we will remove unsupported claims regarding prognostic prediction and future disease progression, ensuring that our conclusions accurately reflect the scope of the analyses performed.

1. Aisen, P. S., et al. (2015). Alzheimer's Disease Neuroimaging Initiative 2 Clinical Core: Progress and plans. Alzheimer's \& Dementia, 11(7), 734--739. \url{https://doi.org/10.1016/j.jalz.2015.05.005}
2. Lin, S.-Y., et al. (2022). The Clinical Course of Early and Late Mild Cognitive Impairment. Frontiers in Neurology, 13, 685636. \url{https://doi.org/10.3389/fneur.2022.685636}
3. Petersen, R. C. (2021). MCI Criteria in ADNI: Meeting Biological Expectations. Neurology, 97(12), 597--599. \url{https://doi.org/10.1212/WNL.0000000000012588}
}

% [Audit Note: Results §2.4 的历史三折分类数值及主张需依据重做的分类实验重新核验；Methods §4.5.2 标题与 R3-3 的方案不同，需归组统一。]

\actiontaken{
We have revised the manuscript accordingly across the following sections:
\begin{itemize}[leftmargin=1.5em, itemsep=2pt]
  \item \textbf{Revised Manuscript}
  \item \textbf{Abstract, Line 60, 68}
  \item \textbf{Introduction, Line 77}
  \item \textbf{Introduction, Line 81}
  \item \textbf{Results: Precise Diagnostic Classification and MCI Stage Discrimination, Line 151----158}
  \item \textbf{Results: Precise Diagnostic Classification and MCI Stage Discrimination, Line 167}
  \item \textbf{Methods: Clinical Cohorts and Diagnostic Subclasses, Line 295}
  \item \textbf{Methods: MCI Stage Discrimination Framework, Line 338----340}
\end{itemize}
}

\changedtext{Revised Manuscript}{%
Manuscript](\url{https://app.notion.com/p/af295d43b2dd836ebbe201bd6d6d1ce0})%
}

\changedtext{Abstract, Line 60, 68}{%
This work advances a scalable, data-efficient tool for the early diagnosis of AD and discrimination between MCI stages.
Conclusion%
}

\changedtext{Abstract, Line 60, 68}{%
Crucially, combining these synthesized PET images with MRI enables precise discrimination between early and late stages of mild cognitive impairment (MCI), advancing a scalable, data-efficient tool for the early diagnosis of AD.%
}

\changedtext{Introduction, Line 77}{%
Despite these advances, a critical limitation remains: the subject-specific precision necessary for personalized diagnosis and MCI stage discrimination is still lacking, which is particularly crucial for a heterogeneous condition like AD\~\\cite.%
}

\changedtext{Introduction, Line 81}{%
When combined with the original MRI scans to train deep learning classifiers\~\\cite, the synthesized PET data enable accurate personalized diagnosis and, most importantly, the differentiation between early and late MCI stages.%
}

\changedtext{Results: Precise Diagnostic Classification and MCI Stage Discrimination, Line 151----158}{%
`Diagnostic Classification and MCI Stage Discrimination`%
}

% [Revision Note: 第一段第一句]
\changedtext{Results: Precise Diagnostic Classification and MCI Stage Discrimination, Line 151----158}{%
To evaluate the clinical utility of the synthetic images for personalized diagnosis and MCI stage discrimination, we conducted classification tasks on the ADNI cohort using a 3-fold cross-validation scheme.%
}

\changedtext{Results: Precise Diagnostic Classification and MCI Stage Discrimination, Line 151----158}{%
Crucially, the combination of baseline MRI and synthetic multi-tracer PET also enabled accurate MCI stage discrimination, achieving 81.96\% accuracy in differentiating EMCI from LMCI, compared to 77.65\% using MRI alone (Fig. 5C).%
}

\changedtext{Results: Precise Diagnostic Classification and MCI Stage Discrimination, Line 151----158}{%
This result highlights the potential of DIReCT++ to distinguish different stages of MCI.%
}

% [Revision Note: LaTeX 第 167 行]
\changedtext{Results: Precise Diagnostic Classification and MCI Stage Discrimination, Line 167}{%
Taken together, these findings confirm the translational potential of DIReCT++, as its synthetic outputs serve as effective proxies for real PET biomarkers in downstream diagnostic tasks, including MCI stage discrimination.%
}

\changedtext{Revised Manuscript}{%
<table header-row=``true''>
<colgroup>
<col />
<col width=``428'' />
</colgroup>
<tr>
<td>
</td>
<td>
</td>
</tr>
<tr>
<td>
`enable precise prognostic stratification of MCI`
</td>
<td>
enable precise MCI stage discrimination
</td>
</tr>
<tr>
<td>
`accurate diagnosis and prognostic stratification of MCI`
</td>
<td>
accurate diagnosis and discrimination between MCI stages
</td>
</tr>
</table>
Importantly, the EMCI-versus-LMCI classification was based on cross-sectional diagnostic categories rather than longitudinal conversion outcomes. Although previous studies have reported differences in subsequent dementia-conversion risk between EMCI and LMCI populations, the present study did not evaluate individual-level disease progression or future conversion to AD. Therefore, the classification results should not be interpreted as evidence of prognostic prediction. Longitudinal validation will be necessary to investigate the prognostic relevance of the synthesized PET images.%
}

% [Revision Note: 采用以下定向修改：]
% [Revision Note: 原稿措辞]
% [Revision Note: 建议替换]
% [Revision Note: 新增limitations]
\changedtext{Methods: Clinical Cohorts and Diagnostic Subclasses, Line 295}{%
To account for the heterogeneity of MCI, we further subclassified individuals based on delayed recall performance on the WMS-R Logical Memory II Story A task: Early MCI (EMCI; n = 291) denoted a milder impairment stage, whereas Late MCI (LMCI; n = 152) denoted a more severe impairment stage.%
}

% [Revision Note: 下一句]
\changedtext{Methods: Clinical Cohorts and Diagnostic Subclasses, Line 295}{%
The remaining MCI participants not classified as EMCI or LMCI were grouped as MCI-unclassified (MCI-U; n = 520).%
}

\changedtext{Methods: MCI Stage Discrimination Framework, Line 338----340}{%
To rigorously evaluate the clinical utility of synthesized PET (Syn-PET) images, we assessed their capacity to support diagnostic classification and MCI stage discrimination across the AD continuum.%
}

% [Revision Note: 第二句]
\changedtext{Methods: MCI Stage Discrimination Framework, Line 338----340}{%
Specifically, we trained identical DenseNet classifiers to discriminate between three clinically salient diagnostic pairs: (1) AD versus CN individuals, a canonical diagnostic boundary; (2) MCI versus CN, representing early disease detection; and (3) EMCI versus late MCI (LMCI), representing different degrees of memory impairment.%
}


\subsection*{R2-5: Information Leakage via Clinical Prompt Conditioning}
\revcomment{R2-5}{5. Potential clinical-information leakage: Age, sex, weight, and cognitive scores are embedded in the text conditioning. Because these variables are strongly associated with diagnosis, ablation experiments excluding cognitive scores are necessary to demonstrate that Syn-PET is not simply re-encoding diagnostic metadata.}

\response{
We thank Reviewer~2 for this insightful and constructive observation.
% [TODO: Insert detailed response for R2-5]
}

\actiontaken{
% [TODO: Insert summary of changes for R2-5]
}

\changedtext{Section X.X, Page Y}{%
% [TODO: Insert revised excerpt for R2-5]
}







\subsection*{R2-6: Subject-Level Isolation and Longitudinal Visit Separation}
\revcomment{R2-6}{> Subject-level splitting must be explicit: Because ADNI contains longitudinal scans, the authors should confirm that all visits from a given participant remain within a single fold to avoid subject leakage between training and testing.}

\response{
We thank the reviewer for highlighting the importance of subject-level data splitting and the potential risk of data leakage arising from longitudinal visits.

Although the original ADNI database contains longitudinal imaging data, the subset used in our study was selected such that only one imaging record was retained per subject. Consequently, no scans from different visits of the same subject could be assigned to different cross-validation folds, thereby preventing subject-level data leakage arising from repeated visits.

We will clarify this data selection criterion in the revised manuscript to avoid ambiguity.
}


\actiontaken{
We have revised the manuscript accordingly across the following sections:
\begin{itemize}[leftmargin=1.5em, itemsep=2pt]
  \item \textbf{Methods: Materials and Data Preprocessing, Line 293}
\end{itemize}
}

\changedtext{Methods: Materials and Data Preprocessing, Line 293}{%
This study leveraged neuroimaging data from 1,857 participants enrolled in the Alzheimer's Disease Neuroimaging Initiative (ADNI), spanning cohorts from ADNI-1 through ADNI-4. All participants underwent structural T1-weighted MRI, while subsets also had $^{18}\text{F}$-FDG-PET ($n = 1,598$) and $^{18}\text{F}$-AV45-PET ($n = 1,365$) scans available, enabling multi-tracer PET synthesis. Only one imaging record per participant was retained for analysis, excluding repeated longitudinal visits from the same participant.%
}


\subsection*{R2-7: Scope of External Validation (OASIS-3 AV45-Only)}
\revcomment{R2-7}{> External validation is incomplete for multi-tracer claims: OASIS-3 external evaluation includes AV45-PET but excludes FDG-PET. Therefore, external generalizability is demonstrated only partially and the multi-tracer generalization claim should be qualified.}

\response{
We thank the reviewer for highlighting the need to clarify the scope of our external validation.

In this study, both AV45-PET and FDG-PET synthesis were evaluated internally on the ADNI dataset, whereas external validation on the independent OASIS-3 dataset was conducted only for AV45-PET. As described in the original manuscript (Section 4.3), FDG-PET scans in OASIS-3 were not included in the external evaluation due to the lack of a standardized PUP processing pipeline.

We agree that the current OASIS-3 results provide evidence of cross-dataset generalizability for AV45-PET synthesis only and cannot be directly extended to FDG-PET synthesis. Accordingly, we will clarify the scope of external validation and refine the relevant statements regarding multi-tracer generalizability in the revised manuscript to ensure that our conclusions accurately reflect the evaluations performed.
}


\actiontaken{
We have revised the manuscript accordingly across the following sections:
\begin{itemize}[leftmargin=1.5em, itemsep=2pt]
  \item \textbf{Abstract: Results, Line 66}
  \item \textbf{Introduction: Summary of Contributions, Line 81}
  \item \textbf{Results: Cross-Dataset Generalizability on OASIS-3, Line 106}
  \item \textbf{Figure 2 Caption, Line 112----114}
\end{itemize}
}

\changedtext{Abstract: Results, Line 66}{%
Extensive evaluations on multi-center datasets demonstrate that DIReCT++ not only produces synthetic PET images ($^{18}\text{F}$-AV-45 and $^{18}\text{F}$-FDG) of superior fidelity but also accurately recapitulates disease-specific patterns, with cross-dataset generalizability demonstrated for $^{18}\text{F}$-AV-45 PET synthesis.%
}

\changedtext{Introduction: Summary of Contributions, Line 81}{%
Comprehensive experiments demonstrate that our model generates multi-tracer PET images with superior fidelity compared to state-of-the-art methods, with cross-dataset generalizability further demonstrated for $^{18}\text{F}$-AV-45 PET synthesis.%
}

\changedtext{Results: Cross-Dataset Generalizability on OASIS-3, Line 106}{%
External validation on the OASIS cohort confirmed that the advantages of DIReCT++ in $^{18}\text{F}$-AV-45 PET synthesis generalize beyond the training data (Fig. 2A).%
}

% [Revision Note: 只增加 in ¹⁸F-AV-45 PET synthesis，其余保持不变。]
\changedtext{Results: Cross-Dataset Generalizability on OASIS-3, Line 106}{%
This robust performance across diverse datasets affirms the state-of-the-art status of DIReCT++ especially for $^{18}\text{F}$-AV-45 PET synthesis, highlighting its superior fidelity and generalization capacity.%
}

\changedtext{Figure 2 Caption, Line 112----114}{%
(A) Radar plots show normalized quantitative results (SSIM, MSE, MAE, PSNR) for FDG-PET and AV45-PET reconstruction on ADNI, and cross-dataset generalization of AV45-PET on OASIS.%
}

\changedtext{Figure 2 Caption, Line 112----114}{%
(C) Generalization to OASIS for AV45-PET using models trained on ADNI.%
}

% [Revision Note: 两处均为局部补充示踪剂名称，不改变图注的其他描述，也不需要修改 Figure 2 的图片内容。]
% [Revision Note: 修改]

\subsection*{R2-8: Non-Significance vs. Formal Equivalence Testing}
\revcomment{R2-8}{> Non-significance does not prove equivalence: Using $p > 0.05$ between real and synthetic PET as evidence of high fidelity is statistically insufficient. Equivalence testing (e.g., TOST) with clinically meaningful equivalence margins is more appropriate.}

\response{
We thank the reviewer for pointing out that the absence of a statistically significant difference does not constitute evidence of equivalence. We agree that relying solely on paired comparisons with $p > 0.05$ between real and synthetic PET images is statistically insufficient to support claims of high fidelity.

To provide a more comprehensive evaluation of regional uptake agreement, we performed additional statistical analyses using the originally preserved subject-level paired data. In addition to paired $t$-tests and multiple-comparison correction, we calculated the mean paired differences and their 95\% confidence intervals, mean absolute error (MAE), Lin's concordance correlation coefficient (CCC), intraclass correlation coefficient (ICC), and Bland--Altman bias and limits of agreement. Together, these metrics provide a more comprehensive assessment of mean bias, regional errors, and subject-level agreement.

The additional analyses showed that, across the 98 brain regions evaluated for each tracer, the median regional CCC values were 0.9745 and 0.9450 for AV45 and FDG, respectively, while the corresponding median regional MAE values were 0.0160 and 0.0159. Under the predefined multiple-comparison correction procedure, no statistically significant mean paired differences were detected in the vast majority of regions and diagnostic groups.

Regarding the suggested two one-sided tests (TOST), we acknowledge the value of formal equivalence testing. However, the regional uptake values in our study were derived from per-image min--max intensity normalization rather than clinically standardized quantitative measurements such as SUV or SUVR. Currently, a well-established equivalence margin with a clear clinical interpretation is lacking for measurements on this normalized intensity scale. Defining an arbitrary equivalence margin without sufficient justification could lead to conclusions that lack clinical interpretability. Therefore, we did not perform TOST to claim clinical equivalence.

Accordingly, we will revise the statistical interpretation in the manuscript to clearly distinguish the absence of statistically significant mean differences, regional uptake agreement, and formal equivalence. We will also moderate the original statements that inferred high fidelity from $p > 0.05$. Furthermore, we will explicitly acknowledge in the Limitations section that our current analyses support agreement in normalized regional uptake values but do not establish clinical quantitative equivalence. Future work will explore PET quantitative evaluation scales with clear biological and clinical interpretations, the definition of appropriate equivalence margins, and corresponding formal equivalence testing methods.
}


\actiontaken{
We have revised the manuscript accordingly across the following sections:
\begin{itemize}[leftmargin=1.5em, itemsep=2pt]
  \item \textbf{Results Section 2.3: Regional Uptake Value Definition, Line 119}
  \item \textbf{Results Section 2.3: Objectives of Regional Statistical Analyses, Line 134}
  \item \textbf{Results Section 2.3: Real-Syn Regional Uptake Concordance, Line 136}
  \item \textbf{Methods Section 4.5.1: Real-Syn Pairwise Consistency Testing, Line 331}
  \item \textbf{Methods Section 4.5.1: AD/CN Between-Group Effect Sizes, Line 333}
  \item \textbf{Methods Section 4.5.1: False Discovery Rate (BH-FDR) Comparison Families, Line 336}
  \item \textbf{Discussion, Line 186}
\end{itemize}
}

\changedtext{Results Section 2.3: Regional Uptake Value Definition, Line 119}{%
To assess the clinical validity of the synthesized PET images, we parcellated both synthetic and real PET scans into anatomical regions using \texttt{mri\_synthseg} (FreeSurfer v7.4.1)~\cite{fischl2012freesurfer,dalca2022synthseg}, a deep learning-based segmentation tool that provides high-resolution cortical and subcortical labeling consistent with the Desikan--Killiany and \textit{aseg} atlases. This approach yielded a total of 98 distinct regions of interest (ROIs), including cortical, subcortical, cerebellar, and brainstem structures. Regional normalized uptake values (arbitrary units) were computed as the mean voxel intensity within each ROI from PET images normalized using per-image min--max intensity scaling. Fig.~\ref{fig:3} illustrates this analysis for key ROIs in two representative ADNI subjects, demonstrating strong regional consistency between synthetic and real images, as well as clear distinctions between syntheses from different diagnostic groups.%
}

% [Revision Note: G4 合并讨论方案，尚待最终确认。]
% [Revision Note: **Original**]
% [Revision Note: ```txt]
% [Revision Note: To assess the clinical validity of the synthesized PET images, we parcellated both synthetic and real PET scans into anatomical regions using \texttt{mri\_synthseg} (FreeSurfer v7.4.1)~\cite{fischl2012freesurfer,dalca2022synthseg}, a deep learning-based segmentation tool that provides high-resolution cortical and subcortical labeling consistent with the Desikan--Killiany and \textit{aseg} atlases. This approach yielded a total of 98 distinct regions of interest (ROIs), including cortical, subcortical, cerebellar, and brainstem structures. Regional standardized uptake value  (SUV) were computed by extracting the mean tracer uptake within each ROI and normalizing to the cerebellar gray matter (for \textsuperscript{18}F-AV-45 PET) or the pons (for \textsuperscript{18}F-FDG PET). Fig.~\ref{fig:3} illustrates this analysis for key ROIs in two representative ADNI subjects, demonstrating strong regional consistency between synthetic and real images, as well as clear distinctions between syntheses from different diagnostic groups.]
% [Revision Note: ```]
% [Revision Note: **Revised**]
% [Revision Note: ```txt]
% [Revision Note: To assess the clinical validity of the synthesized PET images, we parcellated both synthetic and real PET scans into anatomical regions using \texttt{mri\_synthseg} (FreeSurfer v7.4.1)~\cite{fischl2012freesurfer,dalca2022synthseg}, a deep learning-based segmentation tool that provides high-resolution cortical and subcortical labeling consistent with the Desikan--Killiany and \textit{aseg} atlases. This approach yielded a total of 98 distinct regions of interest (ROIs), including cortical, subcortical, cerebellar, and brainstem structures. Regional normalized uptake values (arbitrary units) were computed as the mean voxel intensity within each ROI from PET images normalized using per-image min--max intensity scaling. Fig.~\ref{fig:3} illustrates this analysis for key ROIs in two representative ADNI subjects, demonstrating strong regional consistency between synthetic and real images, as well as clear distinctions between syntheses from different diagnostic groups.]
% [Revision Note: ```]
\changedtext{Results Section 2.3: Objectives of Regional Statistical Analyses, Line 134}{%
For a comprehensive quantitative evaluation, we performed two complementary regional statistical analyses on the ADNI cohort. First, we compared synthetic and real images to assess the normalized regional uptake agreement of DIReCT$++$. Second, we compared the AD and CN groups separately in real and synthetic PET to assess whether disease-related regional differences were retained.%
}

% [Revision Note: G4 合并讨论方案，尚待最终确认。]
% [Revision Note: **Original**]
% [Revision Note: ```txt]
% [Revision Note: For a comprehensive quantitative evaluation, we performed two group-wise statistical comparisons on the ADNI cohort. First, we compared synthetic and real images to validate the regional SUV precision of DIReCT$++$. Second, we compared synthetic images across diagnostic groups (cognitively normal, i.e., CN, MCI, AD) to verify their ability to capture disease-relevant patterns.]
% [Revision Note: ```]
% [Revision Note: **Revised**]
% [Revision Note: ```txt]
% [Revision Note: For a comprehensive quantitative evaluation, we performed two complementary regional statistical analyses on the ADNI cohort. First, we compared synthetic and real images to assess the normalized regional uptake agreement of DIReCT$++$. Second, we compared the AD and CN groups separately in real and synthetic PET to assess whether disease-related regional differences were retained.]
% [Revision Note: ```]
\changedtext{Results Section 2.3: Real-Syn Regional Uptake Concordance, Line 136}{%
Paired-sample $t$-tests between synthetic and real images revealed no statistically significant differences in normalized regional uptake across the principal brain regions shown in Fig.~\ref{fig:4}A, B for the overall cohorts of both tracers ($p>0.05$). Across all 98 ROIs, no overall comparisons remained significant after Benjamini--Hochberg false discovery rate (FDR) correction (all $q>0.05$). Among the diagnosis-stratified paired comparisons, only the right pericalcarine cortex (ROI 2021) in the $^{18}$F-FDG MCI subgroup showed a significant mean paired difference after FDR correction ($q=0.0481$). The median concordance correlation coefficients (CCC) across the 98 ROIs were 0.9745 for $^{18}$F-AV-45 and 0.9450 for $^{18}$F-FDG, with corresponding median mean absolute errors (MAE) of 0.0160 and 0.0159 (normalized intensity units). For the 16 representative tracer--ROI combinations in Fig.~\ref{fig:4}A, B, descriptive regression slopes ranged from 0.9794 to 1.0127 (median, 1.0007) for $^{18}$F-AV-45 and from 0.9844 to 1.0101 (median, 1.0001) for $^{18}$F-FDG, without a consistent direction of deviation from unity. These results indicate close agreement in normalized regional uptake between synthetic and real PET, with detailed statistics for the representative ROIs provided in Supplementary Table~\ref{tab:regional_agreement}.%
}

% [Revision Note: G4 合并讨论方案，尚待最终确认。]
% [Revision Note: **Original**]
% [Revision Note: ```txt]
% [Revision Note: Paired-sample $t$-tests between synthetic and real images revealed no statistically significant differences in mean SUV across all principal brain regions for the CN, MCI, and AD groups ($p>0.05$; Fig.~\ref{fig:4}A, B; figs.~\ref{fig:supp_av45_1}, \ref{fig:supp_av45_2}, \ref{fig:supp_fdg_1}, and \ref{fig:supp_fdg_2}). This high fidelity was consistent for both $^{18}$F-AV-45 (Fig.~\ref{fig:4}A) and $^{18}$F-FDG (Fig.~\ref{fig:4}B), indicating that DIReCT$++$ effectively captures the regional characteristics of the ground truth.]
% [Revision Note: ```]
% [Revision Note: **Revised**]
% [Revision Note: ```txt]
% [Revision Note: Paired-sample $t$-tests between synthetic and real images revealed no statistically significant differences in normalized regional uptake across the principal brain regions shown in Fig.~\ref{fig:4}A, B for the overall cohorts of both tracers ($p>0.05$). Across all 98 ROIs, no overall comparisons remained significant after Benjamini--Hochberg false discovery rate (FDR) correction (all $q>0.05$). Among the diagnosis-stratified paired comparisons, only the right pericalcarine cortex (ROI 2021) in the $^{18}$F-FDG MCI subgroup showed a significant mean paired difference after FDR correction ($q=0.0481$). The median concordance correlation coefficients (CCC) across the 98 ROIs were 0.9745 for $^{18}$F-AV-45 and 0.9450 for $^{18}$F-FDG, with corresponding median mean absolute errors (MAE) of 0.0160 and 0.0159 (normalized intensity units). For the 16 representative tracer--ROI combinations in Fig.~\ref{fig:4}A, B, descriptive regression slopes ranged from 0.9794 to 1.0127 (median, 1.0007) for $^{18}$F-AV-45 and from 0.9844 to 1.0101 (median, 1.0001) for $^{18}$F-FDG, without a consistent direction of deviation from unity. These results indicate close agreement in normalized regional uptake between synthetic and real PET, with detailed statistics for the representative ROIs provided in Supplementary Table~\ref{tab:regional_agreement}.]
% [Revision Note: ```]
\changedtext{Methods Section 4.5.1: Real-Syn Pairwise Consistency Testing, Line 331}{%
\begin{itemize}[leftmargin=1.5em, itemsep=2pt]
\item \textbf{Intra-group consistency}: We performed paired two-sided $t$-tests to compare normalized regional uptake between real and Syn-PET images within each ROI, both across the overall cohorts and within individual diagnostic groups. The null hypothesis was that the mean paired difference in normalized regional uptake was zero. Mean paired differences with 95\% confidence intervals, mean absolute error (MAE), Lin's concordance correlation coefficient (CCC), the two-way absolute-agreement single-measure intraclass correlation coefficient (ICC(A,1)), and Bland--Altman bias and 95\% limits of agreement were also calculated to assess regional agreement. For the 16 representative tracer--ROI combinations shown in Fig.~\ref{fig:4}A, B, we additionally performed subject-level ordinary least-squares regression of synthetic on real normalized regional uptake in the overall cohorts, reporting slopes, intercepts, and their unadjusted 95\% confidence intervals.
\end{itemize}%
}

% [Revision Note: G4 合并讨论方案，尚待最终确认。]
% [Revision Note: **Original**]
% [Revision Note: ```txt]
% [Revision Note: \item \textbf{Intra-group consistency}: For each diagnostic group (AD and CN), we performed paired two-sided $t$-tests to compare the mean standardized uptake values (SUVs) between real and Syn-PET images within each key ROI. The null hypothesis was that the mean difference in SUV between real and synthesized scans is zero. A non-significant result ($p > 0.05$, corrected for multiple comparisons) would indicate high fidelity of Syn-PET to ground-truth PET at the group level.]
% [Revision Note: ```]
% [Revision Note: **Revised**]
% [Revision Note: ```txt]
% [Revision Note: \item \textbf{Intra-group consistency}: We performed paired two-sided $t$-tests to compare normalized regional uptake between real and Syn-PET images within each ROI, both across the overall cohorts and within individual diagnostic groups. The null hypothesis was that the mean paired difference in normalized regional uptake was zero. Mean paired differences with 95\% confidence intervals, mean absolute error (MAE), Lin's concordance correlation coefficient (CCC), the two-way absolute-agreement single-measure intraclass correlation coefficient (ICC(A,1)), and Bland--Altman bias and 95\% limits of agreement were also calculated to assess regional agreement. For the 16 representative tracer--ROI combinations shown in Fig.~\ref{fig:4}A, B, we additionally performed subject-level ordinary least-squares regression of synthetic on real normalized regional uptake in the overall cohorts, reporting slopes, intercepts, and their unadjusted 95\% confidence intervals.]
% [Revision Note: ```]
\changedtext{Methods Section 4.5.1: AD/CN Between-Group Effect Sizes, Line 333}{%
\begin{itemize}[leftmargin=1.5em, itemsep=2pt]
\item \textbf{Inter-group discriminability}: To assess whether Syn-PET preserves diagnostic signal, we conducted independent two-sided Student's $t$-tests (assuming equal variances) comparing mean normalized regional uptake between the AD and CN groups in the same ROIs---first using real PET and then repeating the analysis with Syn-PET. Cohen's $d$ was calculated separately for real and synthetic PET as the AD--CN mean difference divided by the pooled within-group sample standard deviation, to characterize the direction and magnitude of the group differences.
\end{itemize}%
}

% [Revision Note: G4 合并讨论方案，尚待最终确认。]
% [Revision Note: **Original**]
% [Revision Note: ```txt]
% [Revision Note: \item \textbf{Inter-group discriminability}: To assess whether Syn-PET preserves diagnostic signal, we conducted independent two-sided $t$-tests comparing mean SUVs between the AD and CN groups in the same ROIs—first using real PET and then repeating the analysis with Syn-PET. Preservation of statistically significant group differences (e.g., hypometabolism in FDG, elevated amyloid in AV45) in Syn-PET data would demonstrate its capacity to recapitulate clinically meaningful biomarkers.]
% [Revision Note: ```]
% [Revision Note: **Revised**]
% [Revision Note: ```txt]
% [Revision Note: \item \textbf{Inter-group discriminability}: To assess whether Syn-PET preserves diagnostic signal, we conducted independent two-sided Student's $t$-tests (assuming equal variances) comparing mean normalized regional uptake between the AD and CN groups in the same ROIs—first using real PET and then repeating the analysis with Syn-PET. Cohen's $d$ was calculated separately for real and synthetic PET as the AD--CN mean difference divided by the pooled within-group sample standard deviation, to characterize the direction and magnitude of the group differences.]
% [Revision Note: ```]
\changedtext{Methods Section 4.5.1: False Discovery Rate (BH-FDR) Comparison Families, Line 336}{%
All $t$-tests were two-tailed, and $p$-values were adjusted for multiple ROI comparisons using the Benjamini--Hochberg false discovery rate (FDR) procedure~\cite{benjamini1995controlling}. For the paired real--synthetic comparisons, FDR correction was applied across 98 ROIs separately for each tracer and analysis group (overall, AD, CN, EMCI, LMCI, and MCI), yielding 12 families (2 tracers $\times$ 6 analysis groups) and 1,176 tests in total. For the AD--CN comparisons, FDR correction was applied across 98 ROIs separately for each tracer and PET type (real or synthetic), yielding four families and 392 tests in total. Statistical significance was defined as an FDR-adjusted $q<0.05$.%
}

% [Revision Note: G4 合并讨论方案，尚待最终确认。]
% [Revision Note: **Original**]
% [Revision Note: ```txt]
% [Revision Note: All statistical tests were two-tailed, and $p$-values were adjusted for multiple ROI comparisons using the false discovery rate (FDR) procedure~\cite{benjamini1995controlling}. This dual-validation framework ensures that Syn-PET not only approximates real PET quantitatively but also retains its discriminative power for early diagnostic stratification.]
% [Revision Note: ```]
% [Revision Note: **Revised**]
% [Revision Note: ```txt]
% [Revision Note: All $t$-tests were two-tailed, and $p$-values were adjusted for multiple ROI comparisons using the Benjamini--Hochberg false discovery rate (FDR) procedure~\cite{benjamini1995controlling}. For the paired real--synthetic comparisons, FDR correction was applied across 98 ROIs separately for each tracer and analysis group (overall, AD, CN, EMCI, LMCI, and MCI), yielding 12 families (2 tracers $\times$ 6 analysis groups) and 1,176 tests in total. For the AD--CN comparisons, FDR correction was applied across 98 ROIs separately for each tracer and PET type (real or synthetic), yielding four families and 392 tests in total. Statistical significance was defined as an FDR-adjusted $q<0.05$.]
% [Revision Note: ```]
\changedtext{Discussion, Line 186}{%
Several limitations of this study present clear avenues for future work. First, the current text guidance relies on fundamental patient information and imaging knowledge. Incorporating additional clinically accessible data, such as blood biomarkers~\cite{hansson2023blood}, could further enhance subject-level precision, necessitating the development of targeted cross-modal pretraining or prompt-learning strategies. Second, while DIReCT$++$ synthesizes $^{18}$F-AV-45 and $^{18}$F-FDG, extending the framework to include tau-PET imaging would provide a more comprehensive biomarker profile for sensitive screening at the preclinical stage. Third, to maximize clinical utility, future iterations could integrate downstream task guidance (e.g., diagnosis or prognosis labels) directly into the generative model, enabling controllable synthesis~\cite{wang2025flexibly,chang2025controllable} tailored for specific clinical decisions. Regional uptake agreement was evaluated using per-image min--max-normalized intensities rather than standardized SUV or SUVR, and clinically interpretable equivalence margins have not yet been established for these measurements. Future work will explore clinically meaningful PET quantification and justified equivalence margins to enable formal equivalence testing. Finally, the evaluation was conducted on well-characterized but limited research cohorts (ADNI, OASIS). Prospective validation on more diverse, real-world clinical populations (including data from thick-slice or low-field MRI~\cite{sorby2024portable}) is essential to fully establish the generalizability and robustness of the framework.%
}

% [Revision Note: G4 合并讨论方案，尚待最终确认。]
% [Revision Note: **Original**]
% [Revision Note: ```txt]
% [Revision Note: Several limitations of this study present clear avenues for future work. First, the current text guidance relies on fundamental patient information and imaging knowledge. Incorporating additional clinically accessible data, such as blood biomarkers~\cite{hansson2023blood}, could further enhance subject-level precision, necessitating the development of targeted cross-modal pretraining or prompt-learning strategies. Second, while DIReCT$++$ synthesizes $^{18}$F-AV-45 and $^{18}$F-FDG, extending the framework to include tau-PET imaging would provide a more comprehensive biomarker profile for sensitive screening at the preclinical stage. Third, to maximize clinical utility, future iterations could integrate downstream task guidance (e.g., diagnosis or prognosis labels) directly into the generative model, enabling controllable synthesis~\cite{wang2025flexibly,chang2025controllable} tailored for specific clinical decisions. Finally, the evaluation was conducted on well-characterized but limited research cohorts (ADNI, OASIS). Prospective validation on more diverse, real-world clinical populations (including data from thick-slice or low-field MRI~\cite{sorby2024portable}) is essential to fully establish the generalizability and robustness of the framework.]
% [Revision Note: ```]
% [Revision Note: **Revised**]
% [Revision Note: ```txt]
% [Revision Note: Several limitations of this study present clear avenues for future work. First, the current text guidance relies on fundamental patient information and imaging knowledge. Incorporating additional clinically accessible data, such as blood biomarkers~\cite{hansson2023blood}, could further enhance subject-level precision, necessitating the development of targeted cross-modal pretraining or prompt-learning strategies. Second, while DIReCT$++$ synthesizes $^{18}$F-AV-45 and $^{18}$F-FDG, extending the framework to include tau-PET imaging would provide a more comprehensive biomarker profile for sensitive screening at the preclinical stage. Third, to maximize clinical utility, future iterations could integrate downstream task guidance (e.g., diagnosis or prognosis labels) directly into the generative model, enabling controllable synthesis~\cite{wang2025flexibly,chang2025controllable} tailored for specific clinical decisions. Regional uptake agreement was evaluated using per-image min--max-normalized intensities rather than standardized SUV or SUVR, and clinically interpretable equivalence margins have not yet been established for these measurements. Future work will explore clinically meaningful PET quantification and justified equivalence margins to enable formal equivalence testing. Finally, the evaluation was conducted on well-characterized but limited research cohorts (ADNI, OASIS). Prospective validation on more diverse, real-world clinical populations (including data from thick-slice or low-field MRI~\cite{sorby2024portable}) is essential to fully establish the generalizability and robustness of the framework.]
% [Revision Note: ```]
% [Revision Note: #### 补充材料：Sx — 区域一致性与描述性回归（完整 16 行）]
% [Revision Note: - 位置：Supplementary Materials]
% [Revision Note: - TeX 行号：]
% [Revision Note: - 状态：拟新增；编号及最终内容待确认。]
% [Revision Note: Regional uptake agreement and descriptive regression between synthetic and real PET]
% [Revision Note: | Tracer | ROI | N | Mean diff. (95% CI) | MAE | CCC | ICC(A,1) | BA bias (95% LoA) | p | q | OLS slope (95% CI) | OLS intercept (95% CI) |]
% [Revision Note: | --- | --- | --- | --- | --- | --- | --- | --- | --- | --- | --- | --- |]
% [Revision Note: | AV45 | L precuneus (1025) | 1365 | +0.00052 \[−0.00056, +0.00159\] | 0.0163 | 0.9863 | 0.9863 | +0.00052 \[−0.0391, +0.0401\] | 0.3445 | 0.9786 | 1.0127 \[1.0038, 1.0215\] | -0.0064 \[-0.0114, -0.0014\] |]
% [Revision Note: | AV45 | R precuneus (2025) | 1365 | +0.00021 \[−0.00079, +0.00120\] | 0.0150 | 0.9887 | 0.9887 | +0.00021 \[−0.0366, +0.0370\] | 0.6846 | 0.9786 | 1.0008 \[0.9928, 1.0088\] | -0.0002 \[-0.0047, +0.0043\] |]
% [Revision Note: | AV45 | L frontal pole (1032) | 1365 | +0.00033 \[−0.00072, +0.00138\] | 0.0157 | 0.9808 | 0.9808 | +0.00033 \[−0.0385, +0.0392\] | 0.5384 | 0.9786 | 1.0037 \[0.9932, 1.0143\] | -0.0012 \[-0.0055, +0.0032\] |]
% [Revision Note: | AV45 | R frontal pole (2032) | 1364 | +0.00039 \[−0.00068, +0.00147\] | 0.0160 | 0.9815 | 0.9815 | +0.00039 \[−0.0392, +0.0400\] | 0.4746 | 0.9786 | 1.0007 \[0.9904, 1.0110\] | +0.0001 \[-0.0042, +0.0045\] |]
% [Revision Note: | AV45 | L temporal pole (1033) | 1365 | −0.00009 \[−0.00117, +0.00099\] | 0.0162 | 0.9591 | 0.9591 | −0.00009 \[−0.0399, +0.0398\] | 0.8769 | 0.9786 | 0.9794 \[0.9640, 0.9947\] | +0.0078 \[+0.0019, +0.0137\] |]
% [Revision Note: | AV45 | R temporal pole (2033) | 1365 | +0.00031 \[−0.00076, +0.00138\] | 0.0163 | 0.9644 | 0.9644 | +0.00031 \[−0.0392, +0.0399\] | 0.5667 | 0.9786 | 1.0087 \[0.9942, 1.0232\] | -0.0032 \[-0.0092, +0.0027\] |]
% [Revision Note: | AV45 | L superior frontal (1028) | 1365 | −0.00037 \[−0.00146, +0.00073\] | 0.0165 | 0.9806 | 0.9807 | −0.00037 \[−0.0407, +0.0400\] | 0.5102 | 0.9786 | 1.0002 \[0.9896, 1.0107\] | -0.0005 \[-0.0058, +0.0049\] |]
% [Revision Note: | AV45 | R superior frontal (2028) | 1365 | +0.00063 \[−0.00041, +0.00168\] | 0.0157 | 0.9822 | 0.9822 | +0.00063 \[−0.0380, +0.0393\] | 0.2354 | 0.9786 | 0.9970 \[0.9869, 1.0071\] | +0.0021 \[-0.0030, +0.0073\] |]
% [Revision Note: | FDG | L precuneus (1025) | 1598 | −0.00003 \[−0.00100, +0.00093\] | 0.0156 | 0.9573 | 0.9573 | −0.00003 \[−0.0386, +0.0385\] | 0.9433 | 0.9774 | 0.9844 \[0.9699, 0.9990\] | +0.0108 \[+0.0006, +0.0209\] |]
% [Revision Note: | FDG | R precuneus (2025) | 1598 | +0.00010 \[−0.00086, +0.00105\] | 0.0156 | 0.9598 | 0.9598 | +0.00010 \[−0.0380, +0.0382\] | 0.8428 | 0.9774 | 0.9983 \[0.9841, 1.0125\] | +0.0013 \[-0.0086, +0.0111\] |]
% [Revision Note: | FDG | L frontal pole (1032) | 1598 | +0.00024 \[−0.00074, +0.00122\] | 0.0160 | 0.9586 | 0.9586 | +0.00024 \[−0.0389, +0.0394\] | 0.6273 | 0.9410 | 1.0039 \[0.9895, 1.0184\] | -0.0019 \[-0.0100, +0.0061\] |]
% [Revision Note: | FDG | R frontal pole (2032) | 1597 | −0.00032 \[−0.00130, +0.00065\] | 0.0159 | 0.9634 | 0.9634 | −0.00032 \[−0.0394, +0.0387\] | 0.5160 | 0.9410 | 1.0041 \[0.9906, 1.0177\] | -0.0026 \[-0.0103, +0.0050\] |]
% [Revision Note: | FDG | L temporal pole (1033) | 1598 | −0.00034 \[−0.00132, +0.00063\] | 0.0159 | 0.9430 | 0.9430 | −0.00034 \[−0.0392, +0.0385\] | 0.4901 | 0.9410 | 0.9934 \[0.9764, 1.0105\] | +0.0023 \[-0.0046, +0.0091\] |]
% [Revision Note: | FDG | R temporal pole (2033) | 1598 | +0.00068 \[−0.00030, +0.00166\] | 0.0161 | 0.9449 | 0.9449 | +0.00068 \[−0.0386, +0.0400\] | 0.1754 | 0.9410 | 0.9902 \[0.9736, 1.0069\] | +0.0047 \[-0.0022, +0.0116\] |]
% [Revision Note: | FDG | L superior frontal (1028) | 1598 | +0.00093 \[−0.00006, +0.00191\] | 0.0159 | 0.9438 | 0.9438 | +0.00093 \[−0.0385, +0.0403\] | 0.0657 | 0.9410 | 1.0101 \[0.9931, 1.0271\] | -0.0052 \[-0.0156, +0.0051\] |]
% [Revision Note: | FDG | R superior frontal (2028) | 1598 | +0.00081 \[−0.00017, +0.00178\] | 0.0158 | 0.9465 | 0.9465 | +0.00081 \[−0.0382, +0.0398\] | 0.1055 | 0.9410 | 1.0020 \[0.9855, 1.0185\] | -0.0004 \[-0.0104, +0.0096\] |]
% [Revision Note: Supplementary Table note（建议）]
% [Revision Note: Regional uptake and error measurements are expressed in arbitrary units after per-image min--max normalization. Mean paired differences and Bland--Altman bias are defined as synthetic minus real PET. The reported BH-FDR-adjusted q-values correspond to 98 ROI comparisons within each tracer-specific overall cohort. OLS regression was performed with real regional uptake as the independent variable and synthetic regional uptake as the dependent variable. Regression slopes and intercepts are accompanied by unadjusted, descriptive 95% confidence intervals. These regression results do not establish clinical SUV or SUVR calibration.]

\subsection*{R2-9: FDR Multiplicity Correction Families and Reporting}
\revcomment{R2-9}{> Multiple-comparison reporting should be clearer: FDR correction is appropriate, but the statistical family, number of tests, and adjusted p/q-values should be reported explicitly.}

\response{
We thank the reviewer for the valuable suggestion regarding the transparency of multiple-comparison correction reporting. We agree that, although the original manuscript stated that false discovery rate (FDR) correction was applied, the statistical families, number of tests, and adjusted p/q-values were not sufficiently specified.

To address this concern, we have confirmed the original regional statistical data and performed a systematic re-evaluation of the regional comparisons between real and synthetic PET images.

For the paired regional uptake agreement analysis between real and synthetic PET, we performed two-sided paired $t$-tests across all 98 regions of interest (ROIs) for each tracer, both in the overall cohorts and within individual diagnostic subgroups. Benjamini--Hochberg FDR correction was applied separately within each statistical family, comprising 12 families and 1,176 regional comparisons in total. The corresponding results for representative brain regions in the overall cohorts, including raw $p$-values and FDR-adjusted q-values, have been compiled in a new supplementary table.

For the AD versus CN comparisons shown in Figure 4C--F, we additionally performed two-sided independent-samples $t$-tests separately for AV45 Real, AV45 Syn, FDG Real, and FDG Syn PET. Benjamini--Hochberg FDR correction was applied across all 98 ROIs within each of these four statistical families, comprising 392 comparisons in total.

The re-evaluation showed that all 20 significant annotations in Figure 4C/D remained statistically significant after FDR correction, while all 20 non-significant annotations in Figure 4E/F remained non-significant. Sensitivity analyses using Welch's t-test did not change the significance classifications of these regions. Thus, the existing significant and non-significant classifications in Figure 4C--F are supported by the FDR-adjusted results and do not require changes.

In the revised manuscript, we will explicitly report the statistical families and corresponding numbers of tests, clarify the FDR correction procedure and significance thresholds in the Statistical Analysis section and the relevant descriptions of Figure 4, and provide adjusted q-values for representative brain regions to improve the transparency and reproducibility of the statistical analysis.
}


\actiontaken{
We have revised the manuscript accordingly across the following sections:
\begin{itemize}[leftmargin=1.5em, itemsep=2pt]
  \item \textbf{Results Section 2.3: AD/CN Diagnostic Differentiation and Effect Sizes, Line 138}
  \item \textbf{Methods Section 4.5.1: AD/CN Between-Group Effect Sizes, Line 333}
  \item \textbf{Methods Section 4.5.1: False Discovery Rate (BH-FDR) Comparison Families, Line 336}
  \item \textbf{Figure 4 Caption, Line 143----146}
\end{itemize}
}

\changedtext{Results Section 2.3: AD/CN Diagnostic Differentiation and Effect Sizes, Line 138}{%
Critically, independent-sample $t$-tests demonstrated that the significant differences between diagnostic groups observed in real PET data were also observed in the synthetic images (Fig.~\ref{fig:4}C, D; figs.~\ref{fig:supp_disease_av45} and \ref{fig:supp_disease_fdg}). For instance, DIReCT$++$ recapitulated the statistically significant differences between AD and CN groups in the left precuneus for $^{18}$F-AV-45 (real: $q=7.73\times10^{-66}$; synthetic: $q=1.25\times10^{-63}$; Fig.~\ref{fig:4}C) and in the left hippocampus for $^{18}$F-FDG (real: $q=4.27\times10^{-15}$; synthetic: $q=1.27\times10^{-13}$; Fig.~\ref{fig:4}D). Similarly, regional differences that were non-significant in real PET data remained non-significant in the synthetic images after FDR correction (all $q\geq0.05$ for the regions shown in Fig.~\ref{fig:4}E, F). The AD--CN effect-size directions were consistent between real and synthetic PET for all ten regions shown in Fig.~\ref{fig:4}C, D, although the absolute Cohen's $d$ decreased in nine regions (maximum reduction, 0.0771). These findings support the preservation of group-level disease-related regional uptake patterns, while indicating variations in effect-size magnitude. Complete effect-size estimates and FDR-adjusted results for the displayed regions are provided in Supplementary Table~\ref{tab:regional_effect_sizes}.%
}

% [Revision Note: G4 合并讨论方案，尚待最终确认。]
% [Revision Note: **Original**]
% [Revision Note: ```txt]
% [Revision Note: Critically, independent-sample $t$-tests demonstrated that the significant differences between diagnostic groups observed in real PET data were precisely preserved in the synthetic images (Fig.~\ref{fig:4}C, D; figs.~\ref{fig:supp_disease_av45} and \ref{fig:supp_disease_fdg}). For instance, DIReCT$++$ recapitulated the statistically significant differences ($p<0.001$) between AD and CN groups in the precuneus for $^{18}$F-AV-45 (Fig.~\ref{fig:4}C) and in the hippocampus for $^{18}$F-FDG (Fig.~\ref{fig:4}D). Similarly, regional differences that were non-significant in real PET data were also consistently non-significant in the synthetic images ($p>0.05$, Fig.~\ref{fig:4}E, F). These findings confirm that DIReCT$++$ reliably preserves the essential biomarker information necessary for clinical differentiation.]
% [Revision Note: ```]
% [Revision Note: **Revised**]
% [Revision Note: ```txt]
% [Revision Note: Critically, independent-sample $t$-tests demonstrated that the significant differences between diagnostic groups observed in real PET data were also observed in the synthetic images (Fig.~\ref{fig:4}C, D; figs.~\ref{fig:supp_disease_av45} and \ref{fig:supp_disease_fdg}). For instance, DIReCT$++$ recapitulated the statistically significant differences between AD and CN groups in the left precuneus for $^{18}$F-AV-45 (real: $q=7.73\times10^{-66}$; synthetic: $q=1.25\times10^{-63}$; Fig.~\ref{fig:4}C) and in the left hippocampus for $^{18}$F-FDG (real: $q=4.27\times10^{-15}$; synthetic: $q=1.27\times10^{-13}$; Fig.~\ref{fig:4}D). Similarly, regional differences that were non-significant in real PET data remained non-significant in the synthetic images after FDR correction (all $q\geq0.05$ for the regions shown in Fig.~\ref{fig:4}E, F). The AD--CN effect-size directions were consistent between real and synthetic PET for all ten regions shown in Fig.~\ref{fig:4}C, D, although the absolute Cohen's $d$ decreased in nine regions (maximum reduction, 0.0771). These findings support the preservation of group-level disease-related regional uptake patterns, while indicating variations in effect-size magnitude. Complete effect-size estimates and FDR-adjusted results for the displayed regions are provided in Supplementary Table~\ref{tab:regional_effect_sizes}.]
% [Revision Note: ```]
\changedtext{Methods Section 4.5.1: AD/CN Between-Group Effect Sizes, Line 333}{%
\begin{itemize}[leftmargin=1.5em, itemsep=2pt]
\item \textbf{Inter-group discriminability}: To assess whether Syn-PET preserves diagnostic signal, we conducted independent two-sided Student's $t$-tests (assuming equal variances) comparing mean normalized regional uptake between the AD and CN groups in the same ROIs---first using real PET and then repeating the analysis with Syn-PET. Cohen's $d$ was calculated separately for real and synthetic PET as the AD--CN mean difference divided by the pooled within-group sample standard deviation, to characterize the direction and magnitude of the group differences.
\end{itemize}%
}

% [Revision Note: G4 合并讨论方案，尚待最终确认。]
% [Revision Note: **Original**]
% [Revision Note: ```txt]
% [Revision Note: \item \textbf{Inter-group discriminability}: To assess whether Syn-PET preserves diagnostic signal, we conducted independent two-sided $t$-tests comparing mean SUVs between the AD and CN groups in the same ROIs—first using real PET and then repeating the analysis with Syn-PET. Preservation of statistically significant group differences (e.g., hypometabolism in FDG, elevated amyloid in AV45) in Syn-PET data would demonstrate its capacity to recapitulate clinically meaningful biomarkers.]
% [Revision Note: ```]
% [Revision Note: **Revised**]
% [Revision Note: ```txt]
% [Revision Note: \item \textbf{Inter-group discriminability}: To assess whether Syn-PET preserves diagnostic signal, we conducted independent two-sided Student's $t$-tests (assuming equal variances) comparing mean normalized regional uptake between the AD and CN groups in the same ROIs—first using real PET and then repeating the analysis with Syn-PET. Cohen's $d$ was calculated separately for real and synthetic PET as the AD--CN mean difference divided by the pooled within-group sample standard deviation, to characterize the direction and magnitude of the group differences.]
% [Revision Note: ```]
\changedtext{Methods Section 4.5.1: False Discovery Rate (BH-FDR) Comparison Families, Line 336}{%
All $t$-tests were two-tailed, and $p$-values were adjusted for multiple ROI comparisons using the Benjamini--Hochberg false discovery rate (FDR) procedure~\cite{benjamini1995controlling}. For the paired real--synthetic comparisons, FDR correction was applied across 98 ROIs separately for each tracer and analysis group (overall, AD, CN, EMCI, LMCI, and MCI), yielding 12 families (2 tracers $\times$ 6 analysis groups) and 1,176 tests in total. For the AD--CN comparisons, FDR correction was applied across 98 ROIs separately for each tracer and PET type (real or synthetic), yielding four families and 392 tests in total. Statistical significance was defined as an FDR-adjusted $q<0.05$.%
}

% [Revision Note: G4 合并讨论方案，尚待最终确认。]
% [Revision Note: **Original**]
% [Revision Note: ```txt]
% [Revision Note: All statistical tests were two-tailed, and $p$-values were adjusted for multiple ROI comparisons using the false discovery rate (FDR) procedure~\cite{benjamini1995controlling}. This dual-validation framework ensures that Syn-PET not only approximates real PET quantitatively but also retains its discriminative power for early diagnostic stratification.]
% [Revision Note: ```]
% [Revision Note: **Revised**]
% [Revision Note: ```txt]
% [Revision Note: All $t$-tests were two-tailed, and $p$-values were adjusted for multiple ROI comparisons using the Benjamini--Hochberg false discovery rate (FDR) procedure~\cite{benjamini1995controlling}. For the paired real--synthetic comparisons, FDR correction was applied across 98 ROIs separately for each tracer and analysis group (overall, AD, CN, EMCI, LMCI, and MCI), yielding 12 families (2 tracers $\times$ 6 analysis groups) and 1,176 tests in total. For the AD--CN comparisons, FDR correction was applied across 98 ROIs separately for each tracer and PET type (real or synthetic), yielding four families and 392 tests in total. Statistical significance was defined as an FDR-adjusted $q<0.05$.]
% [Revision Note: ```]
\changedtext{Figure 4 Caption, Line 143----146}{%
\textbf{Caption:} 
\textbf{Regional comparison of mean normalized uptake between real and MRI-synthetic PET images.}
(A,B) Violin plots show regional mean normalized uptake (arbitrary units) across cortical and subcortical regions for real and synthetic PET, comparing AV45 (A) and FDG (B). Statistical significance was assessed for each region. (C--F) AD--CN regional comparisons. Violin plots display mean normalized uptake in AD (red) and CN (green) subjects for selected regions. Each region uses paired half-violins: left (solid) = real PET, right (dashed) = synthetic PET; colors match across modalities. Solid lines indicate real PET significance, dashed lines indicate synthetic PET significance. The statistical annotations in panels (C--F) are consistent with the revised BH-FDR analysis: \texttt{****} corresponds to $q<10^{-4}$ and \texttt{ns} to $q\geq0.05$. (C,D) show representative regions with significant AD--CN differences for AV45 and FDG, whereas (E,F) show selected regions without significant AD--CN differences after FDR correction. Cblm-WM denotes Cerebellum-White-Matter. Additional regional agreement and effect-size analyses are provided in the Supplementary Materials.
\\[4pt]%
}

% [Revision Note: G4 合并讨论方案，尚待最终确认。]
% [Revision Note: **Original**]
% [Revision Note: ```txt]
% [Revision Note: \caption{]
% [Revision Note: \textbf{Regional comparison of mean uptake between real and MRI-synthetic PET images.}]
% [Revision Note: (A,B) Violin plots show regional mean uptake across cortical and subcortical regions for real and synthetic PET, comparing AV45 (A) and FDG (B). Statistical significance was assessed for each region. (C–F) Disease-related patterns. Violin plots display mean uptake in AD (red) and CN (green) subjects for selected regions. Each region uses paired half-violins: left (solid) = real PET, right (dashed) = synthetic PET; colors match across modalities. Solid lines indicate real PET significance, dashed lines indicate synthetic PET significance. (C,D) show representative disease-related regions for AV45 and FDG, (E,F) show non–disease-related regions. Cblm-WM denotes Cerebellum-White-Matter. Additional analyses are in the Supplementary Materials.]
% [Revision Note: }]
% [Revision Note: ```]
% [Revision Note: **Revised**]
% [Revision Note: ```txt]
% [Revision Note: \caption{]
% [Revision Note: \textbf{Regional comparison of mean normalized uptake between real and MRI-synthetic PET images.}]
% [Revision Note: (A,B) Violin plots show regional mean normalized uptake (arbitrary units) across cortical and subcortical regions for real and synthetic PET, comparing AV45 (A) and FDG (B). Statistical significance was assessed for each region. (C--F) AD--CN regional comparisons. Violin plots display mean normalized uptake in AD (red) and CN (green) subjects for selected regions. Each region uses paired half-violins: left (solid) = real PET, right (dashed) = synthetic PET; colors match across modalities. Solid lines indicate real PET significance, dashed lines indicate synthetic PET significance. The statistical annotations in panels (C--F) are consistent with the revised BH-FDR analysis: \texttt{****} corresponds to $q<10^{-4}$ and \texttt{ns} to $q\geq0.05$. (C,D) show representative regions with significant AD--CN differences for AV45 and FDG, whereas (E,F) show selected regions without significant AD--CN differences after FDR correction. Cblm-WM denotes Cerebellum-White-Matter. Additional regional agreement and effect-size analyses are provided in the Supplementary Materials.]
% [Revision Note: }]
% [Revision Note: ```]
% [Revision Note: #### 补充材料：Sy — AD/CN 效应量与 FDR（20 行）]
% [Revision Note: - 位置：Supplementary Materials]
% [Revision Note: - TeX 行号：]
% [Revision Note: - 状态：拟新增；编号及最终内容待确认。]
% [Revision Note: AD/CN effect sizes and FDR-adjusted regional comparisons]
% [Revision Note: R2-9 / R3-5 · Figure 4C–F]
% [Revision Note: 20 行]
% [Revision Note: 新增一张表，报告 Real/Syn 的 AD/CN Cohen's d，以及同一区域各自的 BH-FDR q 值。]
% [Revision Note: 第二张表建议采用以下精简结构，数值来自已核验的 R3-5 CSV。]
% [Revision Note: | Panel / Tracer | ROI ID | Real d | Syn d | Real q | Syn q |]
% [Revision Note: | --- | --- | --- | --- | --- | --- |]
% [Revision Note: | C / AV45 | 1025 | 1.6783 | 1.6437 | 7.73e−66 | 1.25e−63 |]
% [Revision Note: | C / AV45 | 1009 | 1.6831 | 1.6337 | 4.68e−66 | 4.74e−63 |]
% [Revision Note: | C / AV45 | 1015 | 1.6482 | 1.6236 | 4.67e−64 | 1.72e−62 |]
% [Revision Note: | C / AV45 | 1030 | 1.5780 | 1.5434 | 5.64e−60 | 8.51e−58 |]
% [Revision Note: | C / AV45 | 1034 | 1.5133 | 1.5288 | 2.84e−56 | 5.61e−57 |]
% [Revision Note: | D / FDG | 1010 | −1.2531 | −1.1774 | 2.69e−48 | 2.10e−43 |]
% [Revision Note: | D / FDG | 2010 | −1.3972 | −1.3201 | 1.18e−57 | 1.94e−52 |]
% [Revision Note: | D / FDG | 1031 | −0.4848 | −0.4308 | 2.14e−9 | 1.10e−7 |]
% [Revision Note: | D / FDG | 17 | −0.6401 | −0.6042 | 4.27e−15 | 1.27e−13 |]
% [Revision Note: | D / FDG | 53 | −0.7518 | −0.7189 | 4.59e−20 | 1.54e−18 |]
% [Revision Note: | E / AV45 | 28 | −0.1050 | −0.0933 | 0.2368 | 0.2858 |]
% [Revision Note: | E / AV45 | 60 | −0.1211 | −0.0986 | 0.1729 | 0.2618 |]
% [Revision Note: | E / AV45 | 53 | −0.0068 | 0.0006 | 0.9373 | 0.9949 |]
% [Revision Note: | E / AV45 | 7 | −0.0745 | −0.1010 | 0.3996 | 0.2529 |]
% [Revision Note: | E / AV45 | 17 | 0.0701 | 0.1232 | 0.4241 | 0.1667 |]
% [Revision Note: | F / FDG | 10 | −0.0097 | 0.0582 | 0.9008 | 0.4667 |]
% [Revision Note: | F / FDG | 49 | −0.0668 | −0.0530 | 0.3964 | 0.4991 |]
% [Revision Note: | F / FDG | 1019 | −0.1294 | −0.0913 | 0.1035 | 0.2506 |]
% [Revision Note: | F / FDG | 1011 | −0.1239 | −0.0978 | 0.1171 | 0.2229 |]
% [Revision Note: | F / FDG | 2013 | −0.0493 | −0.0435 | 0.5303 | 0.5744 |]
% [Revision Note: 正式表格还应包含解剖名称与 AD/CN 样本数；具体信息均可直接从 [完整 Cohen's d CSV](https://github.com/linshuijin6/DIRECTPP/blob/main/reports/r3_5_effect_size/figure4_effect_sizes.csv?utm_source=chatgpt.com) 生成，不需要人工重新抄录数据。]
% [Revision Note: 新增表注建议]
% [Revision Note: ```txt]
% [Revision Note: Cohen's d was calculated separately for real and synthetic PET as the AD--CN mean difference divided by the pooled within-group sample standard deviation. Positive values indicate higher normalized regional uptake in AD than CN. The reported q-values were obtained from the revised two-sided Student's t-test analysis with Benjamini--Hochberg FDR correction across 98 ROIs within each tracer and PET type. All uptake measurements were derived from per-image min--max-normalized PET images and do not represent SUV or SUVR. Differences in Cohen's d do not establish clinical equivalence or non-equivalence.]
% [Revision Note: ```]

\subsection*{R2-10: Subject-Level Agreement, Calibration, and Regional Concordance}
\revcomment{R2-10}{> PET agreement metrics should be strengthened: SSIM, PSNR, MSE, and MAE are useful image-similarity metrics but are insufficient for biomarker agreement. ICC/CCC, Bland-Altman analysis, regional SUV bias, and calibration analyses would improve clinical validation.}

\response{
We thank the reviewer for the valuable suggestion to strengthen the evaluation of PET biomarker agreement. We agree that image-similarity metrics such as SSIM, PSNR, MSE, and MAE alone are insufficient to characterize subject-level regional uptake agreement and systematic differences between synthetic and real PET images.

To address this concern, we performed additional analyses using the original subject-level paired data across all 98 regions of interest (ROIs) for each tracer. Specifically, we calculated Lin's concordance correlation coefficient (CCC), the intraclass correlation coefficient (ICC(A,1)), mean paired differences with 95\% confidence intervals, and Bland--Altman bias and 95\% limits of agreement. These complementary metrics provide a more comprehensive assessment of regional errors, subject-level agreement, and systematic differences. Across the 98 ROIs, the median regional CCC values were 0.9745 for AV45 and 0.9450 for FDG.

Furthermore, to examine the relationship between synthetic and real regional uptake values, we conducted subject-level ordinary least-squares regression analyses for the 16 representative tracer--ROI combinations shown in Figure 4A/B. Real PET regional uptake was used as the independent variable and synthetic PET regional uptake as the dependent variable. Regression slopes, intercepts, and their corresponding 95\% confidence intervals were calculated.

The regression slopes ranged from 0.9794 to 1.0127 (median: 1.0007) for AV45 and from 0.9844 to 1.0101 (median: 1.0001) for FDG. These representative regions did not exhibit a consistent pattern of linear compression or expansion, although modest deviations from a unit slope were observed in some regions. The regression results provide additional information on how synthetic regional uptake varies with real uptake levels and complement the CCC, ICC, and Bland--Altman analyses.

Importantly, the regional uptake measurements in this study were derived from per-image min--max normalized PET intensities rather than clinically standardized quantitative SUV or SUVR measurements. Therefore, the reported agreement and regression results characterize normalized regional uptake only and should not be interpreted as evidence of clinical SUV/SUVR agreement, absolute quantitative accuracy, or clinically validated calibration. We will explicitly clarify this distinction and its implications in the revised manuscript. Future work will investigate PET synthesis and validation approaches that preserve clinically meaningful quantitative uptake measurements.

In the revised manuscript, we will expand the regional agreement supplementary table to include regression slopes, intercepts, and their 95\% confidence intervals for the 16 representative tracer--ROI combinations, alongside the existing agreement and bias metrics. We will also clarify the corresponding statistical methods and limitations.
}


\actiontaken{
We have revised the manuscript accordingly across the following sections:
\begin{itemize}[leftmargin=1.5em, itemsep=2pt]
  \item \textbf{Results Section 2.3: Regional Uptake Value Definition, Line 119}
  \item \textbf{Results Section 2.3: Real-Syn Regional Uptake Concordance, Line 136}
  \item \textbf{Methods Section 4.5.1: Real-Syn Pairwise Consistency Testing, Line 331}
  \item \textbf{Discussion, Line 186}
\end{itemize}
}

\changedtext{Results Section 2.3: Regional Uptake Value Definition, Line 119}{%
To assess the clinical validity of the synthesized PET images, we parcellated both synthetic and real PET scans into anatomical regions using \texttt{mri\_synthseg} (FreeSurfer v7.4.1)~\cite{fischl2012freesurfer,dalca2022synthseg}, a deep learning-based segmentation tool that provides high-resolution cortical and subcortical labeling consistent with the Desikan--Killiany and \textit{aseg} atlases. This approach yielded a total of 98 distinct regions of interest (ROIs), including cortical, subcortical, cerebellar, and brainstem structures. Regional normalized uptake values (arbitrary units) were computed as the mean voxel intensity within each ROI from PET images normalized using per-image min--max intensity scaling. Fig.~\ref{fig:3} illustrates this analysis for key ROIs in two representative ADNI subjects, demonstrating strong regional consistency between synthetic and real images, as well as clear distinctions between syntheses from different diagnostic groups.%
}

% [Revision Note: G4 合并讨论方案，尚待最终确认。]
% [Revision Note: **Original**]
% [Revision Note: ```txt]
% [Revision Note: To assess the clinical validity of the synthesized PET images, we parcellated both synthetic and real PET scans into anatomical regions using \texttt{mri\_synthseg} (FreeSurfer v7.4.1)~\cite{fischl2012freesurfer,dalca2022synthseg}, a deep learning-based segmentation tool that provides high-resolution cortical and subcortical labeling consistent with the Desikan--Killiany and \textit{aseg} atlases. This approach yielded a total of 98 distinct regions of interest (ROIs), including cortical, subcortical, cerebellar, and brainstem structures. Regional standardized uptake value  (SUV) were computed by extracting the mean tracer uptake within each ROI and normalizing to the cerebellar gray matter (for \textsuperscript{18}F-AV-45 PET) or the pons (for \textsuperscript{18}F-FDG PET). Fig.~\ref{fig:3} illustrates this analysis for key ROIs in two representative ADNI subjects, demonstrating strong regional consistency between synthetic and real images, as well as clear distinctions between syntheses from different diagnostic groups.]
% [Revision Note: ```]
% [Revision Note: **Revised**]
% [Revision Note: ```txt]
% [Revision Note: To assess the clinical validity of the synthesized PET images, we parcellated both synthetic and real PET scans into anatomical regions using \texttt{mri\_synthseg} (FreeSurfer v7.4.1)~\cite{fischl2012freesurfer,dalca2022synthseg}, a deep learning-based segmentation tool that provides high-resolution cortical and subcortical labeling consistent with the Desikan--Killiany and \textit{aseg} atlases. This approach yielded a total of 98 distinct regions of interest (ROIs), including cortical, subcortical, cerebellar, and brainstem structures. Regional normalized uptake values (arbitrary units) were computed as the mean voxel intensity within each ROI from PET images normalized using per-image min--max intensity scaling. Fig.~\ref{fig:3} illustrates this analysis for key ROIs in two representative ADNI subjects, demonstrating strong regional consistency between synthetic and real images, as well as clear distinctions between syntheses from different diagnostic groups.]
% [Revision Note: ```]
\changedtext{Results Section 2.3: Real-Syn Regional Uptake Concordance, Line 136}{%
Paired-sample $t$-tests between synthetic and real images revealed no statistically significant differences in normalized regional uptake across the principal brain regions shown in Fig.~\ref{fig:4}A, B for the overall cohorts of both tracers ($p>0.05$). Across all 98 ROIs, no overall comparisons remained significant after Benjamini--Hochberg false discovery rate (FDR) correction (all $q>0.05$). Among the diagnosis-stratified paired comparisons, only the right pericalcarine cortex (ROI 2021) in the $^{18}$F-FDG MCI subgroup showed a significant mean paired difference after FDR correction ($q=0.0481$). The median concordance correlation coefficients (CCC) across the 98 ROIs were 0.9745 for $^{18}$F-AV-45 and 0.9450 for $^{18}$F-FDG, with corresponding median mean absolute errors (MAE) of 0.0160 and 0.0159 (normalized intensity units). For the 16 representative tracer--ROI combinations in Fig.~\ref{fig:4}A, B, descriptive regression slopes ranged from 0.9794 to 1.0127 (median, 1.0007) for $^{18}$F-AV-45 and from 0.9844 to 1.0101 (median, 1.0001) for $^{18}$F-FDG, without a consistent direction of deviation from unity. These results indicate close agreement in normalized regional uptake between synthetic and real PET, with detailed statistics for the representative ROIs provided in Supplementary Table~\ref{tab:regional_agreement}.%
}

% [Revision Note: G4 合并讨论方案，尚待最终确认。]
% [Revision Note: **Original**]
% [Revision Note: ```txt]
% [Revision Note: Paired-sample $t$-tests between synthetic and real images revealed no statistically significant differences in mean SUV across all principal brain regions for the CN, MCI, and AD groups ($p>0.05$; Fig.~\ref{fig:4}A, B; figs.~\ref{fig:supp_av45_1}, \ref{fig:supp_av45_2}, \ref{fig:supp_fdg_1}, and \ref{fig:supp_fdg_2}). This high fidelity was consistent for both $^{18}$F-AV-45 (Fig.~\ref{fig:4}A) and $^{18}$F-FDG (Fig.~\ref{fig:4}B), indicating that DIReCT$++$ effectively captures the regional characteristics of the ground truth.]
% [Revision Note: ```]
% [Revision Note: **Revised**]
% [Revision Note: ```txt]
% [Revision Note: Paired-sample $t$-tests between synthetic and real images revealed no statistically significant differences in normalized regional uptake across the principal brain regions shown in Fig.~\ref{fig:4}A, B for the overall cohorts of both tracers ($p>0.05$). Across all 98 ROIs, no overall comparisons remained significant after Benjamini--Hochberg false discovery rate (FDR) correction (all $q>0.05$). Among the diagnosis-stratified paired comparisons, only the right pericalcarine cortex (ROI 2021) in the $^{18}$F-FDG MCI subgroup showed a significant mean paired difference after FDR correction ($q=0.0481$). The median concordance correlation coefficients (CCC) across the 98 ROIs were 0.9745 for $^{18}$F-AV-45 and 0.9450 for $^{18}$F-FDG, with corresponding median mean absolute errors (MAE) of 0.0160 and 0.0159 (normalized intensity units). For the 16 representative tracer--ROI combinations in Fig.~\ref{fig:4}A, B, descriptive regression slopes ranged from 0.9794 to 1.0127 (median, 1.0007) for $^{18}$F-AV-45 and from 0.9844 to 1.0101 (median, 1.0001) for $^{18}$F-FDG, without a consistent direction of deviation from unity. These results indicate close agreement in normalized regional uptake between synthetic and real PET, with detailed statistics for the representative ROIs provided in Supplementary Table~\ref{tab:regional_agreement}.]
% [Revision Note: ```]
\changedtext{Methods Section 4.5.1: Real-Syn Pairwise Consistency Testing, Line 331}{%
\begin{itemize}[leftmargin=1.5em, itemsep=2pt]
\item \textbf{Intra-group consistency}: We performed paired two-sided $t$-tests to compare normalized regional uptake between real and Syn-PET images within each ROI, both across the overall cohorts and within individual diagnostic groups. The null hypothesis was that the mean paired difference in normalized regional uptake was zero. Mean paired differences with 95\% confidence intervals, mean absolute error (MAE), Lin's concordance correlation coefficient (CCC), the two-way absolute-agreement single-measure intraclass correlation coefficient (ICC(A,1)), and Bland--Altman bias and 95\% limits of agreement were also calculated to assess regional agreement. For the 16 representative tracer--ROI combinations shown in Fig.~\ref{fig:4}A, B, we additionally performed subject-level ordinary least-squares regression of synthetic on real normalized regional uptake in the overall cohorts, reporting slopes, intercepts, and their unadjusted 95\% confidence intervals.
\end{itemize}%
}

% [Revision Note: G4 合并讨论方案，尚待最终确认。]
% [Revision Note: **Original**]
% [Revision Note: ```txt]
% [Revision Note: \item \textbf{Intra-group consistency}: For each diagnostic group (AD and CN), we performed paired two-sided $t$-tests to compare the mean standardized uptake values (SUVs) between real and Syn-PET images within each key ROI. The null hypothesis was that the mean difference in SUV between real and synthesized scans is zero. A non-significant result ($p > 0.05$, corrected for multiple comparisons) would indicate high fidelity of Syn-PET to ground-truth PET at the group level.]
% [Revision Note: ```]
% [Revision Note: **Revised**]
% [Revision Note: ```txt]
% [Revision Note: \item \textbf{Intra-group consistency}: We performed paired two-sided $t$-tests to compare normalized regional uptake between real and Syn-PET images within each ROI, both across the overall cohorts and within individual diagnostic groups. The null hypothesis was that the mean paired difference in normalized regional uptake was zero. Mean paired differences with 95\% confidence intervals, mean absolute error (MAE), Lin's concordance correlation coefficient (CCC), the two-way absolute-agreement single-measure intraclass correlation coefficient (ICC(A,1)), and Bland--Altman bias and 95\% limits of agreement were also calculated to assess regional agreement. For the 16 representative tracer--ROI combinations shown in Fig.~\ref{fig:4}A, B, we additionally performed subject-level ordinary least-squares regression of synthetic on real normalized regional uptake in the overall cohorts, reporting slopes, intercepts, and their unadjusted 95\% confidence intervals.]
% [Revision Note: ```]
\changedtext{Discussion, Line 186}{%
Several limitations of this study present clear avenues for future work. First, the current text guidance relies on fundamental patient information and imaging knowledge. Incorporating additional clinically accessible data, such as blood biomarkers~\cite{hansson2023blood}, could further enhance subject-level precision, necessitating the development of targeted cross-modal pretraining or prompt-learning strategies. Second, while DIReCT$++$ synthesizes $^{18}$F-AV-45 and $^{18}$F-FDG, extending the framework to include tau-PET imaging would provide a more comprehensive biomarker profile for sensitive screening at the preclinical stage. Third, to maximize clinical utility, future iterations could integrate downstream task guidance (e.g., diagnosis or prognosis labels) directly into the generative model, enabling controllable synthesis~\cite{wang2025flexibly,chang2025controllable} tailored for specific clinical decisions. Regional uptake agreement was evaluated using per-image min--max-normalized intensities rather than standardized SUV or SUVR, and clinically interpretable equivalence margins have not yet been established for these measurements. Future work will explore clinically meaningful PET quantification and justified equivalence margins to enable formal equivalence testing. Finally, the evaluation was conducted on well-characterized but limited research cohorts (ADNI, OASIS). Prospective validation on more diverse, real-world clinical populations (including data from thick-slice or low-field MRI~\cite{sorby2024portable}) is essential to fully establish the generalizability and robustness of the framework.%
}

% [Revision Note: G4 合并讨论方案，尚待最终确认。]
% [Revision Note: **Original**]
% [Revision Note: ```txt]
% [Revision Note: Several limitations of this study present clear avenues for future work. First, the current text guidance relies on fundamental patient information and imaging knowledge. Incorporating additional clinically accessible data, such as blood biomarkers~\cite{hansson2023blood}, could further enhance subject-level precision, necessitating the development of targeted cross-modal pretraining or prompt-learning strategies. Second, while DIReCT$++$ synthesizes $^{18}$F-AV-45 and $^{18}$F-FDG, extending the framework to include tau-PET imaging would provide a more comprehensive biomarker profile for sensitive screening at the preclinical stage. Third, to maximize clinical utility, future iterations could integrate downstream task guidance (e.g., diagnosis or prognosis labels) directly into the generative model, enabling controllable synthesis~\cite{wang2025flexibly,chang2025controllable} tailored for specific clinical decisions. Finally, the evaluation was conducted on well-characterized but limited research cohorts (ADNI, OASIS). Prospective validation on more diverse, real-world clinical populations (including data from thick-slice or low-field MRI~\cite{sorby2024portable}) is essential to fully establish the generalizability and robustness of the framework.]
% [Revision Note: ```]
% [Revision Note: **Revised**]
% [Revision Note: ```txt]
% [Revision Note: Several limitations of this study present clear avenues for future work. First, the current text guidance relies on fundamental patient information and imaging knowledge. Incorporating additional clinically accessible data, such as blood biomarkers~\cite{hansson2023blood}, could further enhance subject-level precision, necessitating the development of targeted cross-modal pretraining or prompt-learning strategies. Second, while DIReCT$++$ synthesizes $^{18}$F-AV-45 and $^{18}$F-FDG, extending the framework to include tau-PET imaging would provide a more comprehensive biomarker profile for sensitive screening at the preclinical stage. Third, to maximize clinical utility, future iterations could integrate downstream task guidance (e.g., diagnosis or prognosis labels) directly into the generative model, enabling controllable synthesis~\cite{wang2025flexibly,chang2025controllable} tailored for specific clinical decisions. Regional uptake agreement was evaluated using per-image min--max-normalized intensities rather than standardized SUV or SUVR, and clinically interpretable equivalence margins have not yet been established for these measurements. Future work will explore clinically meaningful PET quantification and justified equivalence margins to enable formal equivalence testing. Finally, the evaluation was conducted on well-characterized but limited research cohorts (ADNI, OASIS). Prospective validation on more diverse, real-world clinical populations (including data from thick-slice or low-field MRI~\cite{sorby2024portable}) is essential to fully establish the generalizability and robustness of the framework.]
% [Revision Note: ```]
% [Revision Note: #### 补充材料：Sx — 区域一致性与描述性回归（完整 16 行）]
% [Revision Note: - 位置：Supplementary Materials]
% [Revision Note: - TeX 行号：]
% [Revision Note: - 状态：拟新增；编号及最终内容待确认。]
% [Revision Note: Regional uptake agreement and descriptive regression between synthetic and real PET]
% [Revision Note: | Tracer | ROI | N | Mean diff. (95% CI) | MAE | CCC | ICC(A,1) | BA bias (95% LoA) | p | q | OLS slope (95% CI) | OLS intercept (95% CI) |]
% [Revision Note: | --- | --- | --- | --- | --- | --- | --- | --- | --- | --- | --- | --- |]
% [Revision Note: | AV45 | L precuneus (1025) | 1365 | +0.00052 \[−0.00056, +0.00159\] | 0.0163 | 0.9863 | 0.9863 | +0.00052 \[−0.0391, +0.0401\] | 0.3445 | 0.9786 | 1.0127 \[1.0038, 1.0215\] | -0.0064 \[-0.0114, -0.0014\] |]
% [Revision Note: | AV45 | R precuneus (2025) | 1365 | +0.00021 \[−0.00079, +0.00120\] | 0.0150 | 0.9887 | 0.9887 | +0.00021 \[−0.0366, +0.0370\] | 0.6846 | 0.9786 | 1.0008 \[0.9928, 1.0088\] | -0.0002 \[-0.0047, +0.0043\] |]
% [Revision Note: | AV45 | L frontal pole (1032) | 1365 | +0.00033 \[−0.00072, +0.00138\] | 0.0157 | 0.9808 | 0.9808 | +0.00033 \[−0.0385, +0.0392\] | 0.5384 | 0.9786 | 1.0037 \[0.9932, 1.0143\] | -0.0012 \[-0.0055, +0.0032\] |]
% [Revision Note: | AV45 | R frontal pole (2032) | 1364 | +0.00039 \[−0.00068, +0.00147\] | 0.0160 | 0.9815 | 0.9815 | +0.00039 \[−0.0392, +0.0400\] | 0.4746 | 0.9786 | 1.0007 \[0.9904, 1.0110\] | +0.0001 \[-0.0042, +0.0045\] |]
% [Revision Note: | AV45 | L temporal pole (1033) | 1365 | −0.00009 \[−0.00117, +0.00099\] | 0.0162 | 0.9591 | 0.9591 | −0.00009 \[−0.0399, +0.0398\] | 0.8769 | 0.9786 | 0.9794 \[0.9640, 0.9947\] | +0.0078 \[+0.0019, +0.0137\] |]
% [Revision Note: | AV45 | R temporal pole (2033) | 1365 | +0.00031 \[−0.00076, +0.00138\] | 0.0163 | 0.9644 | 0.9644 | +0.00031 \[−0.0392, +0.0399\] | 0.5667 | 0.9786 | 1.0087 \[0.9942, 1.0232\] | -0.0032 \[-0.0092, +0.0027\] |]
% [Revision Note: | AV45 | L superior frontal (1028) | 1365 | −0.00037 \[−0.00146, +0.00073\] | 0.0165 | 0.9806 | 0.9807 | −0.00037 \[−0.0407, +0.0400\] | 0.5102 | 0.9786 | 1.0002 \[0.9896, 1.0107\] | -0.0005 \[-0.0058, +0.0049\] |]
% [Revision Note: | AV45 | R superior frontal (2028) | 1365 | +0.00063 \[−0.00041, +0.00168\] | 0.0157 | 0.9822 | 0.9822 | +0.00063 \[−0.0380, +0.0393\] | 0.2354 | 0.9786 | 0.9970 \[0.9869, 1.0071\] | +0.0021 \[-0.0030, +0.0073\] |]
% [Revision Note: | FDG | L precuneus (1025) | 1598 | −0.00003 \[−0.00100, +0.00093\] | 0.0156 | 0.9573 | 0.9573 | −0.00003 \[−0.0386, +0.0385\] | 0.9433 | 0.9774 | 0.9844 \[0.9699, 0.9990\] | +0.0108 \[+0.0006, +0.0209\] |]
% [Revision Note: | FDG | R precuneus (2025) | 1598 | +0.00010 \[−0.00086, +0.00105\] | 0.0156 | 0.9598 | 0.9598 | +0.00010 \[−0.0380, +0.0382\] | 0.8428 | 0.9774 | 0.9983 \[0.9841, 1.0125\] | +0.0013 \[-0.0086, +0.0111\] |]
% [Revision Note: | FDG | L frontal pole (1032) | 1598 | +0.00024 \[−0.00074, +0.00122\] | 0.0160 | 0.9586 | 0.9586 | +0.00024 \[−0.0389, +0.0394\] | 0.6273 | 0.9410 | 1.0039 \[0.9895, 1.0184\] | -0.0019 \[-0.0100, +0.0061\] |]
% [Revision Note: | FDG | R frontal pole (2032) | 1597 | −0.00032 \[−0.00130, +0.00065\] | 0.0159 | 0.9634 | 0.9634 | −0.00032 \[−0.0394, +0.0387\] | 0.5160 | 0.9410 | 1.0041 \[0.9906, 1.0177\] | -0.0026 \[-0.0103, +0.0050\] |]
% [Revision Note: | FDG | L temporal pole (1033) | 1598 | −0.00034 \[−0.00132, +0.00063\] | 0.0159 | 0.9430 | 0.9430 | −0.00034 \[−0.0392, +0.0385\] | 0.4901 | 0.9410 | 0.9934 \[0.9764, 1.0105\] | +0.0023 \[-0.0046, +0.0091\] |]
% [Revision Note: | FDG | R temporal pole (2033) | 1598 | +0.00068 \[−0.00030, +0.00166\] | 0.0161 | 0.9449 | 0.9449 | +0.00068 \[−0.0386, +0.0400\] | 0.1754 | 0.9410 | 0.9902 \[0.9736, 1.0069\] | +0.0047 \[-0.0022, +0.0116\] |]
% [Revision Note: | FDG | L superior frontal (1028) | 1598 | +0.00093 \[−0.00006, +0.00191\] | 0.0159 | 0.9438 | 0.9438 | +0.00093 \[−0.0385, +0.0403\] | 0.0657 | 0.9410 | 1.0101 \[0.9931, 1.0271\] | -0.0052 \[-0.0156, +0.0051\] |]
% [Revision Note: | FDG | R superior frontal (2028) | 1598 | +0.00081 \[−0.00017, +0.00178\] | 0.0158 | 0.9465 | 0.9465 | +0.00081 \[−0.0382, +0.0398\] | 0.1055 | 0.9410 | 1.0020 \[0.9855, 1.0185\] | -0.0004 \[-0.0104, +0.0096\] |]
% [Revision Note: Supplementary Table note（建议）]
% [Revision Note: Regional uptake and error measurements are expressed in arbitrary units after per-image min--max normalization. Mean paired differences and Bland--Altman bias are defined as synthetic minus real PET. The reported BH-FDR-adjusted q-values correspond to 98 ROI comparisons within each tracer-specific overall cohort. OLS regression was performed with real regional uptake as the independent variable and synthetic regional uptake as the dependent variable. Regression slopes and intercepts are accompanied by unadjusted, descriptive 95% confidence intervals. These regression results do not establish clinical SUV or SUVR calibration.]

\subsection*{R2-11: Classification Uncertainty and Confidence Intervals}
\revcomment{R2-11}{11. Classification uncertainty is underreported: ACC, AUC, sensitivity, and specificity should be accompanied by fold-wise mean $\pm$ SD or 95\% confidence intervals rather than isolated point estimates.}

\response{
We thank Reviewer~2 for this insightful and constructive observation.
% [TODO: Insert detailed response for R2-11]
}

\actiontaken{
% [TODO: Insert summary of changes for R2-11]
}

\changedtext{Section X.X, Page Y}{%
% [TODO: Insert revised excerpt for R2-11]
}







\subsection*{R2-12: Statistical Significance Testing of Classification Baselines}
\revcomment{R2-12}{12. Claims of statistical superiority require direct tests: Statements such as ``significantly outperforming'' should be supported by formal paired comparisons, such as DeLong testing for AUC or paired bootstrap/McNemar tests where appropriate.}

\response{
We thank Reviewer~2 for this insightful and constructive observation.
% [TODO: Insert detailed response for R2-12]
}

\actiontaken{
% [TODO: Insert summary of changes for R2-12]
}

\changedtext{Section X.X, Page Y}{%
% [TODO: Insert revised excerpt for R2-12]
}







\subsection*{R2-13: One-Step Sampling Terminology and Rectified Flow Integration}
\revcomment{R2-13}{> One-step sampling terminology requires correction: The phrase ``exact one-step sampling'' is too strong for a distilled end-point predictor. The manuscript should distinguish exact ODE integration from learned one-step approximation.}

\response{
We thank the reviewer for pointing out the imprecision in our description of one-step sampling. We agree that the continuous ODE process underlying rectified flow should be distinguished from the learned one-step approximation employed in DIReCT++.

Specifically, DIReCT++ generates the target PET image directly from the input MRI in a single forward pass, which does not imply exact integration of the underlying ODE. Accordingly, we will revise the expression ``exact one-step sampling'' to more accurately describe the learned one-step generation process.
}


\actiontaken{
We have revised the manuscript accordingly across the following sections:
\begin{itemize}[leftmargin=1.5em, itemsep=2pt]
  \item \textbf{Methods Section 4.2.3: Theoretical Formulation of One-Step Sampling, Line 283}
  \item \textbf{Methods Section 4.2.3: Formulation of Eq. (7), Line 286}
\end{itemize}
}

\changedtext{Methods Section 4.2.3: Theoretical Formulation of One-Step Sampling, Line 283}{%
A key advantage of rectified flow is its potential to support efficient one-step generation through distillation.%
}

\changedtext{Methods Section 4.2.3: Formulation of Eq. (7), Line 286}{%
Mathematically, this learned one-step approximation is expressed as:%
}


\subsection*{R2-14: VLM Guidance and Adapter Implementation Details}
\revcomment{R2-14}{14. Domain-adapted VLM details are insufficient: The manuscript should report embedding dimensionality, PET image-feature extraction, prompt construction, affine-adapter training, threshold $\tau$ selection, and optimization details to ensure reproducibility.}

\response{
We thank Reviewer~2 for this insightful and constructive observation.
% [TODO: Insert detailed response for R2-14]
}

\actiontaken{
% [TODO: Insert summary of changes for R2-14]
}

\changedtext{Section X.X, Page Y}{%
% [TODO: Insert revised excerpt for R2-14]
}







\subsection*{R2-15: Equation 1, Inter-Tracer Similarity vs. Distinctness, and Tau}
\revcomment{R2-15}{> Equation (1) needs stronger justification: The semantic-similarity constraint between FDG and AV45 embeddings may conflict with the stated goal of preserving inter-tracer differences. A sensitivity analysis for $\tau$ and an explicit ablation are recommended.}

\response{
We thank the reviewer for the valuable comments regarding the rationale and parameter sensitivity of the cross-tracer semantic similarity constraint. We agree that its design motivation and its role in balancing text--PET feature alignment and tracer-specific representations warrant further clarification.

In DIReCT++, we employ a pre-trained BiomedCLIP model to encode the textual descriptions of FDG and AV45, followed by two independent affine adapters that align the corresponding textual embeddings with their respective PET image features. However, optimizing only the individual text--PET alignment objectives may cause the two adapted guidance vectors to become highly similar, potentially diminishing the inter-tracer distinctions present in the original textual representations. We therefore introduce a cross-tracer similarity constraint, not to maximize the similarity between the two representations, but to regulate their relative similarity during feature alignment.

Regarding the selection of $\tau$=0.82, we would like to clarify its specific basis. Before affine adaptation, the cosine similarity between the original FDG and AV45 textual embeddings produced by the pre-trained BiomedCLIP model was 0.8205374, which rounds to 0.82. Accordingly, we selected $\tau$=0.82 as the target cross-tracer similarity, using the original relationship between the two textual representations as a reference. This allows the adapters to align the textual embeddings with their corresponding PET image features while maintaining approximately the relative similarity originally encoded by BiomedCLIP. Thus, $\tau$=0.82 was not an arbitrarily assigned value, but was determined from the representational characteristics of the pre-trained text encoder.

We also acknowledge that Equation (1) in the original manuscript presented a simplified inequality constraint to convey the underlying design principle of cross-tracer semantic consistency. In the actual adapter optimization, this principle was implemented using a differentiable squared penalty, `100 × (cos(z\_FDG, z\_AV45) $-$ $\tau$)²`, which encourages the cosine similarity between the adapted tracer embeddings to approach the specified target $\tau$. This formulation penalizes deviations both above and below the target, thereby preventing the two tracer representations from becoming increasingly similar without constraint during text--PET alignment. We will clarify the implemented optimization objective and the rationale for selecting $\tau$ in the revised manuscript.

To further validate this design, we conducted additional adapter-level parameter sensitivity and cross-tracer constraint ablation experiments. Using identical two-fold data splits, original text and PET features, initialization procedures, and optimization settings, we evaluated $\tau$ values of 0.70, 0.82, 0.90, and 0.98, together with an additional configuration in which the cross-tracer similarity penalty was completely removed.

The results demonstrate that $\tau$ effectively regulates the relative similarity between the adapted tracer guidance vectors. As $\tau$ increased from 0.70 to 0.98, the mean cosine similarity between the adapted FDG and AV45 embeddings increased from 0.7032 to 0.9807. At the originally adopted $\tau$=0.82, the resulting similarities were 0.8213 and 0.8215 for the two folds, respectively, closely matching the specified target. In contrast, removing the cross-tracer penalty increased these similarities to 0.9900 and 0.9894, indicating that optimizing text--PET alignment alone caused the two guidance vectors to become highly similar.

Further analysis showed that the individual text--PET feature alignment also increased with $\tau$. At $\tau$=0.82, the mean alignment similarities were 0.9687 for FDG and 0.9661 for AV45. Although removing the cross-tracer constraint further increased these alignment values, it simultaneously brought the cosine similarity between the two adapted tracer embeddings close to 0.99. These findings indicate that maximizing individual text--PET alignment does not necessarily preserve the original representational distinctions between tracers. Instead, the cross-tracer similarity term provides a controllable mechanism for balancing individual feature alignment and inter-tracer representational separation.

Therefore, the rationale for selecting $\tau$=0.82 is to use the original similarity between the BiomedCLIP textual representations as a reference, maintaining substantial text--PET feature alignment while preventing the two adapted guidance vectors from becoming excessively similar.

The additional experiments directly demonstrate the influence of $\tau$ on adapter-level representation learning and the relative geometry of the tracer embeddings. Since these analyses focus on the representation level rather than retraining PET synthesis models under different $\tau$ settings, we present them as supporting evidence for the parameter selection and similarity constraint mechanism, rather than as an optimization study of PET synthesis performance.

We will incorporate the rationale for $\tau$=0.82, the actual optimization formulation, and the additional sensitivity and ablation results into the revised manuscript to improve the methodological clarity and reproducibility of our approach.
}

% [Audit Note: 清单已勾选，但主页面状态仍标为待定；参数及历史实现的表述需作者最终确认。]

\actiontaken{
We have revised the manuscript accordingly across the following sections:
\begin{itemize}[leftmargin=1.5em, itemsep=2pt]
  \item \textbf{Methods: Domain-Adapted VLM Formulation, Line 229----241}
  \item \textbf{Results: The DIReCT++ Framework Overview, Line 88}
  \item \textbf{Supplementary Materials: Adapter Ablation Study}
\end{itemize}
}

\changedtext{Methods: Domain-Adapted VLM Formulation, Line 229----241}{%
where $\alpha_1, \beta_1 $ are learnable scalar weights and $\alpha_2, \beta_2 $ are learnable bias vectors. We choose scalar scaling rather than full linear projection to minimize trainable parameters and avoid overfitting, given the limited size of medical imaging datasets. The optimal parameters are obtained by minimizing the following objective:
\begin{equation}
\begin{aligned}
\mathcal{L}\_{\text{align}} ={}\&
-\frac{1}{N\_f}\sum\_{i=1}^{N\_f}
\operatorname{sim}\big(Adapt\_f(c\_f), f\_{f,i}\big) \\
\&-\frac{1}{N\_a}\sum\_{j=1}^{N\_a}
\operatorname{sim}\big(Adapt\_a(c\_a), f\_{a,j}\big) \\
\&+\lambda\_{\text{cross}}
\left[
\operatorname{sim}\big(Adapt\_f(c\_f), Adapt\_a(c\_a)\big)-\tau
\right]^2,
\end{aligned}
\label{eq:align_opt}
\end{equation}
where $\operatorname{sim}(\cdot,\cdot)$ denotes cosine similarity, $f_{f,i}$ and $f_{a,j}$ denote the normalized PET image features of individual FDG and AV45 training subjects, and $N_f$ and $N_a$ are the corresponding numbers of available features. We set $\lambda_{\text{cross}}=100$ and $\tau=0.82$, the latter corresponding to the cosine similarity of the original FDG and AV45 BiomedCLIP textual embeddings (0.8205) before affine adaptation. The cross-tracer penalty encourages the adapted guidance vectors to retain approximately this reference similarity while being aligned with their respective PET image features, thereby regulating their relative representational separation.

By solving~\eqref{eq:align_opt}, we obtain the modality-specific guidance vectors $Adapt_f(c_f)$ and $Adapt_a(c_a)$, which are subsequently used to steer the rectified flow translation process in a domain-informed manner.%
}

% [Revision Note: - 保留原稿关于两个独立 affine adapters 的结构与解释。]
% [Revision Note: - 将概念性不等式改为实际优化使用的平方惩罚目标。]
% [Revision Note: - 明确 \\(\\tau=0.82\\) 来源于适配前两种文本向量的余弦相似度 0.8205。]
% [Revision Note: - 使图文对齐项与实际的逐受试者余弦均值计算一致。]
% [Revision Note: - 保留原有公式标签 `eq:align_opt`，不影响后续交叉引用。]
% [Revision Note: 这部分不需要在正文额外解释为什么原稿采用简化形式。相关说明已经在 Response 中给出，正文只需准确呈现方法。]
\changedtext{Results: The DIReCT++ Framework Overview, Line 88}{%
An additional adapter-level sensitivity analysis showed that varying $\tau$ from 0.70 to 0.98 changed the mean cross-tracer embedding similarity from 0.7032 to 0.9807. At the adopted $\tau=0.82$, the similarity was 0.8214, compared with 0.9897 when the cross-tracer penalty was removed, demonstrating its role in regulating inter-tracer representational similarity (Table~\ref{tab:supp_tau_adapter}).%
}

% [Revision Note: 只报告关键发现，不在主文重复五组实验的全部指标。新增表格放在 Supplementary Materials 中，供审稿人检查。]
% [Revision Note: 这里刻意使用 adapter-level sensitivity analysis 和 embedding similarity，避免将新增表示层结果误写成 PET 合成性能比较。]
% [Revision Note: 原段落完整保留，在其后新增：]
\changedtext{Supplementary Materials: Adapter Ablation Study}{%
\begin{center}
\textbf{Caption:} \textbf{Adapter-level sensitivity analysis of the cross-tracer similarity constraint.}
Values are descriptive means across two folds. FDG and AV45 alignment denote the mean cosine similarity between each adapted tracer embedding and the corresponding individual PET image features. The cross-tracer similarity is the cosine similarity between the two adapted tracer embeddings. The analysis used 1,856 QC-filtered ADNI subjects. Only training-fold PET features were used for adapter optimization, and held-out features were used solely for evaluation. No cross-term indicates complete removal of the cross-tracer penalty.\\[4pt]
\label{tab:supp_tau_adapter}
\centering
\small
\begin{tabular}{lccccc}
\toprule
\textbf{Setting} &
\multicolumn{2}{c}{\textbf{Training alignment}} &
\multicolumn{2}{c}{\textbf{Held-out alignment}} &
\textbf{Cross-tracer} \\
 & FDG & AV45 & FDG & AV45 & similarity \\
\midrule
$\tau=0.70$ & 0.9437 & 0.9410 & 0.9436 & 0.9410 & 0.7032 \\
$\tau=0.82$ & 0.9687 & 0.9661 & 0.9687 & 0.9661 & 0.8214 \\
$\tau=0.90$ & 0.9838 & 0.9813 & 0.9838 & 0.9812 & 0.9018 \\
$\tau=0.98$ & 0.9948 & 0.9923 & 0.9948 & 0.9923 & 0.9807 \\
No cross-term & 0.9952 & 0.9927 & 0.9952 & 0.9926 & 0.9897 \\
\bottomrule
\end{tabular}
\end{center}%
}

% [Revision Note: 该表包含完整的五组 Adapter 实验，直接对应 R2-15 的参数敏感性与约束消融要求。数据采用两折描述性均值，保留了训练折与留出折的区别。]
% [Revision Note: 你已经生成的 `tau_sensitivity.png` 也可以作为补充图，但我认为目前新增这张表已经能够支撑正文中的结果描述，没有必要为了这一条意见扩大图表修改范围。]
% [Revision Note: 该位置没有对应原文，为新增内容。现有补充材料的末项是软件依赖表，原有内容全部保留。]
% [Revision Note: 建议新增一张表格：]

\subsection*{R2-16: Demographic Confounding and Conditioning Variables}
\revcomment{R2-16}{16. Demographic confounding needs control: ADNI groups differ significantly in age, sex, and other variables that are also used in the conditioning prompt. Matched, adjusted, or subgroup analyses are needed to quantify confounding effects.}

\response{
We thank Reviewer~2 for this insightful and constructive observation.
% [TODO: Insert detailed response for R2-16]
}

\actiontaken{
% [TODO: Insert summary of changes for R2-16]
}

\changedtext{Section X.X, Page Y}{%
% [TODO: Insert revised excerpt for R2-16]
}







\subsection*{R2-17: Baseline Adaptation, Tuning, and Fair Comparison Protocols}
\revcomment{R2-17}{17. Baseline adaptation and fairness should be documented: Some comparator methods originate from segmentation or MRI-to-CT synthesis. The authors should clearly describe how each baseline was adapted and confirm identical data splits, preprocessing, tuning effort, and training budgets.}

\response{
We thank Reviewer~2 for this insightful and constructive observation.
% [TODO: Insert detailed response for R2-17]
}

\actiontaken{
% [TODO: Insert summary of changes for R2-17]
}

\changedtext{Section X.X, Page Y}{%
% [TODO: Insert revised excerpt for R2-17]
}







\subsection*{R2-18: Computational Efficiency, Parameter Counts, and Hardware Benchmarks}
\revcomment{R2-18}{18. Computational-efficiency claims need direct evidence: The manuscript reports approximately 138 GPU-hours for training, but should also report inference time per 3D volume, parameter count, FLOPs, memory use, and direct runtime comparisons with diffusion/flow baselines.}

\response{
We thank Reviewer~2 for this insightful and constructive observation.
% [TODO: Insert detailed response for R2-18]
}

\actiontaken{
% [TODO: Insert summary of changes for R2-18]
}

\changedtext{Section X.X, Page Y}{%
% [TODO: Insert revised excerpt for R2-18]
}







\subsection*{R2-19: Code, Seed, Split, and Model Weight Release}
\revcomment{R2-19}{19. Reproducibility is a strength but can be improved: Public code, trained models, software versions, and hardware details are valuable. The authors should additionally provide fixed data-split files, random seeds, and a repository commit/version used for the reported experiments.}

\response{
We thank Reviewer~2 for this insightful and constructive observation.
% [TODO: Insert detailed response for R2-19]
}

\actiontaken{
% [TODO: Insert summary of changes for R2-19]
}

\changedtext{Section X.X, Page Y}{%
% [TODO: Insert revised excerpt for R2-19]
}







\subsection*{R2-20: Clinical Translation Scope and Translational Boundaries}
\revcomment{R2-20}{> Clinical translation claims should be moderated: Statements regarding widespread screening, clinically actionable biomarkers, and early intervention exceed the evidence from retrospective ADNI/OASIS cohorts. Prospective, diverse, multi-center clinical validation is needed before such claims are justified.}

\response{
We appreciate the reviewer's thoughtful suggestions regarding the clinical translation of our work.

One of the primary motivations of DIReCT++ is to explore the feasibility of synthesizing multi-tracer PET images from routine MRI, potentially alleviating the challenges associated with PET acquisition costs, radiation exposure, and limited accessibility. Our evaluations on the ADNI and OASIS-3 datasets provide preliminary evidence of the framework's PET synthesis performance, while regional uptake analyses and downstream classification on ADNI further support its potential clinical relevance.

Nevertheless, translating these findings from retrospective research cohorts into clinical practice requires more comprehensive evidence. We agree with the reviewer that the current results are insufficient to support the direct application of synthetic PET in routine clinical screening or treatment decision-making. Accordingly, we will moderate statements regarding widespread screening, clinically actionable biomarkers, and early intervention to better distinguish the demonstrated findings from potential future clinical applications. We will also clarify that prospective validation in more diverse, real-world clinical populations remains necessary.
}

% [Audit Note: Discussion 第 177 行旧方案含 “widespread screening roughly”，与弱化临床筛查主张的意图不一致；需在 G7 整合时重审。]

\actiontaken{
We have revised the manuscript accordingly across the following sections:
\begin{itemize}[leftmargin=1.5em, itemsep=2pt]
  \item \textbf{Introduction, Line 81}
  \item \textbf{Discussion, Line 175}
  \item \textbf{Discussion, Line 177}
  \item \textbf{Discussion, Line 179}
  \item \textbf{Discussion, Line 188}
\end{itemize}
}

\changedtext{Introduction, Line 81}{%
Our work proposes a translational approach to precision biomarker synthesis, offering a scalable and data-efficient strategy with potential applications in the early assessment of dementia.%
}

\changedtext{Discussion, Line 175}{%
Beyond its technical merits, the primary significance of DIReCT++ lies in its potential to improve access to biomarker-related imaging information, thereby addressing a major challenge in the early assessment of dementia.%
}

\changedtext{Discussion, Line 175}{%
By generating synthetic estimates of both amyloid and metabolic PET from a single routine MRI scan, our framework offers a scalable approach for multimodal imaging research.%
}

\changedtext{Discussion, Line 175}{%
The potential clinical relevance of this approach is supported by the observed consistency of disease-related regional uptake patterns between synthetic and real PET (Figs. 3 and 4), as well as by our results on personalized MCI diagnosis and stratification (Fig. 5B, C), where the combination of MRI and synthetic multi-tracer PET achieved performance comparable to MRI with real PET and outperformed models using MRI or real PET alone.%
}

\changedtext{Discussion, Line 177}{%
It could enable widespread screening roughly for individuals at the preclinical or prodromal stages in primary care settings, and provide a basis for exploring more accessible assessment strategies.%
}

\changedtext{Discussion, Line 177}{%
DIReCT++ thus offers a potential foundation for more accessible multi-modal imaging approaches in future clinical workflows.%
}

\changedtext{Discussion, Line 179}{%
DIReCT++ moves beyond this by achieving precise generation of subject-specific PET images that recapitulate disease-related regional uptake patterns.%
}

\changedtext{Discussion, Line 188}{%
By generating synthetic PET images that support diagnostic classification and discrimination between MCI stages, our work offers a potential approach to bridging the gap between the biological definition of AD and its clinical application.%
}

% [Revision Note: 保留原来的句式及“连接 AD 生物学研究与临床应用”这一科学意义，但将其定位于潜在应用，而非已经实现的临床诊断能力。]
\changedtext{Discussion, Line 188}{%
DIReCT++ represents a step toward scalable, accessible, and precise tools for early assessment of neurodegenerative disorders.%
}

% [Revision Note: 将 early intervention 调整为 early assessment，使结论与本研究以影像合成和下游分类为核心的验证范围更加一致。]
% [Revision Note: # Reviewer 3 — Major]
